% =============================================================================
% Topological Feature Learning for Parameter Estimation in Neyman-Scott Processes
%
% STRUCTURING PASS 3 -- the paper is rebuilt around one headline model:
% persistence images of the H0 diagram under a single DTM_5 filtration, a
% shared-weight CNN encoder, concatenation with log N, and the per-diagram
% coverage-quantile calibration (q = 0.95). Multi-scale fusion, the H1 channel,
% CoordConv, the encoder/fusion cross and the three alternative vectorizations
% are retained, but as clearly-marked branches in Section 7.3 rather than as
% parts of the method. Prose is spliced from the previous draft wherever
% possible; new text is confined to connective tissue and to claims that had to
% change because the model changed.
%
% HOW TO READ THIS FILE
%   * Every section opens with a % NOTE block: PURPOSE.
%   * \todo and \fix have been cleared from Sections 1-7 and the appendices;
%     the author will re-add them. Sections 8 and 9 are bullet skeletons, not
%     prose.
%   * Surviving annotations are only:
%       \expt -- EXPERIMENT required to support the claim above (orange)
%       \tab  -- TABLE required here       (teal)
%       \fig  -- FIGURE required here      (magenta)
% =============================================================================

\documentclass[11pt,letterpaper]{article}

\usepackage[margin=1in]{geometry}
\usepackage{mathpazo}
\usepackage{microtype}

\usepackage{amsmath}
\usepackage{amssymb}
\usepackage{amsthm}
\usepackage{mathtools}
\usepackage{dsfont}

\usepackage{graphicx}
\graphicspath{{../figs/}}
\usepackage{booktabs}
\usepackage{enumitem}

% ----------------------------------------------------------------------------
% Float and caption handling
%   Captions in this draft are long and doubled as running commentary; small
%   type keeps them from dominating the page.
% ----------------------------------------------------------------------------
\usepackage[font=small,labelfont=bf,skip=6pt]{caption}
\setlength{\tabcolsep}{5pt}
\renewcommand{\arraystretch}{1.08}

\usepackage{natbib}
\usepackage{xcolor}
\usepackage[hidelinks]{hyperref}

% ----------------------------------------------------------------------------
% Macros (unused ones from the AAAI draft have been deleted)
% ----------------------------------------------------------------------------
\newcommand{\R}{\mathbb{R}}
\newcommand{\N}{\mathbb{N}}
\newcommand{\Z}{\mathbb{Z}}
\newcommand{\E}{\mathbb{E}}
\newcommand{\1}{\mathds{1}}
\newcommand{\pc}{\mathbf{x}}
\newcommand{\simpcom}{\Sigma}
\newcommand{\filt}{\mathcal{F}}
\newcommand{\kpar}{\kappa}
\newcommand{\mo}{\mu}

% ----------------------------------------------------------------------------
% Draft annotations
%   \todo -- prose to write            (red)
%   \fix  -- proposed replacement text (blue)
%   \tab  -- TABLE required here       (green)
%   \fig  -- FIGURE required here      (magenta)
%   \expt -- EXPERIMENT required to support the claim above (orange)
%            (named \expt, not \exp, since \exp is LaTeX's built-in
%            exponential-function operator)
% ----------------------------------------------------------------------------
\newif\ifdraft
\drafttrue
\ifdraft
  \newcommand{\todo}[1]{\textcolor{red}{[TODO: #1]}}
  \newcommand{\fix}[1]{\textcolor{blue}{[FIX: #1]}}
  \newcommand{\tab}[1]{\textcolor{teal}{\textbf{[TABLE:] }#1}}
  \newcommand{\fig}[1]{\textcolor{magenta}{\textbf{[FIGURE:] }#1}}
  \newcommand{\expt}[1]{\textcolor{orange}{\textbf{[EXPT:] }#1}}
\else
  \newcommand{\todo}[1]{}\newcommand{\fix}[1]{}
  \newcommand{\tab}[1]{}\newcommand{\fig}[1]{}\newcommand{\expt}[1]{}
\fi

% ----------------------------------------------------------------------------
% Theorem environments (trimmed to those actually used)
% ----------------------------------------------------------------------------
\theoremstyle{plain}
\newtheorem{theorem}{Theorem}[section]
\newtheorem{claim}[theorem]{Claim}

\newtheoremstyle{definitioncolon}%
  {\topsep}{\topsep}{\normalfont}{}{\bfseries}{}{.5em}%
  {\thmname{#1}\thmnumber{ #2}\thmnote{: #3}}
\theoremstyle{definitioncolon}
\newtheorem{definition}[theorem]{Definition}

\newtheorem*{note}{Note}
\newtheorem*{example}{Example}
\newtheorem*{remark}{Remark}

% ----------------------------------------------------------------------------
\title{Topological Feature Learning for Parameter Estimation\\
       in Neyman-Scott Processes}
\author{Flavio Gualtieri}
\date{\today}

\begin{document}
\maketitle


% #############################################################################
% ABSTRACT
% #############################################################################

\begin{abstract}
The parameters of spatial point processes are interpretable scientific
quantities, but the methods available to estimate them rely on closed-form
expressions for hand-picked summary statistics such as the $K$ and $g$
functions. This reliance limits them: a closed form must be re-derived for each
new family of processes, and the choice of statistic discards information about
the point pattern before any estimation takes place. We propose a deep learning
estimator that instead uses vectorized persistence summaries as features, which
require no closed form and no choice of summary statistic. We pursue two
objectives. Firstly, we compare several topological summaries against the
estimation task directly, so that their relative efficacies indicate which
structure each one retains. Secondly, we use the strongest summary to build an
estimator for the Thomas and nested Thomas processes under a
common reparameterization, and evaluate it against classical baselines on shared
test data. Our contributions include a per-diagram calibration scheme for
vectorizing across a heterogeneous population of persistence diagrams, and a
simulator validated against an exact finite-window second-order identity.
On the Thomas process, our estimator reaches a mean loss of
$\mathcal{L}=0.126\pm0.008$, within noise of a specialized classical-summary
neural baseline ($0.115\pm0.007$) and roughly an order of magnitude below
minimum contrast, which is both worse on average ($0.827\pm0.646$ fit to $K$;
$6.167\pm0.331$ fit to $g$) and unstable across realizations. Because it
requires no closed-form summary statistic, the same estimator applies
unchanged to the harder nested Thomas process ($\mathcal{L}=0.196\pm0.006$),
for which no minimum-contrast baseline can be constructed at all.
\end{abstract}


% #############################################################################
% 1. INTRODUCTION
% #############################################################################
% NOTE -- PURPOSE: motivate parameter estimation, state what is wrong with the
% closed-form estimators, and say why persistence features are the response.
% #############################################################################

\section{Introduction}
\label{sec:intro}

% --- (i) motivation ----------------------------------------------------------
The study of spatial point processes finds applications in astronomy, cell
biology, epidemiology, etc REF. Many natural phenomena - such as galaxy clustering REF -
are modelled using SPPs, so one can hope to leverage an understanding of these
processes to obtain novel insights.
% NOTE: pick ONE application and carry it through the whole paper as the running
% example -- galaxy clustering is the natural choice given the related work.

The parameters of spatial point processes are of particular interest because 
they are interpretable scientific quantities. \todo{For example ...}
% NOTE: give an example of a parameter read as a scientific quantity.

% --- (ii) limitations of existing estimators ---------------------------------
A number of parameter estimation methods have been proposed, typically relying
on closed-form expressions for summary statistics, such as the $K$ and $L$
functions (see Section \ref{ssec:spp}). Minimum contrast estimation fits parameters by numerically
minimizing a discrepancy between the empirical and theoretical curve, most
commonly for $K$ or $g$ REF [Diggle 2003]. Palm likelihood instead
builds a composite likelihood from the process's Palm intensity, and has been
developed specifically for Neyman-Scott processes REF [Tanaka 2008].
Fully Bayesian approaches treat the unobserved parent locations as latent
variables and sample the posterior over $\theta$ by MCMC REF
[Moller 2004].

These methods have three limitations we hope to address.

Firstly, they rely on the existence of closed-form expressions of
summary statistics which may not exist for every process. \todo{For example ...}

Secondly, classical parameter estimation methods suffer from 
hyperparameter sensitivity, specifically the minimum contrast method
depends on $(r_{max}, q, p)$. \todo{ParamNet too. Justify}

Finally, classical methods become unstable as the process approaches complete
spatial randomness (near-Poisson regime). As the overlap index $\nu$ grows the clusters merge, so
$\kappa$ and $\mu$ cease to be separately identifiable from a single
realization: a small parent intensity with large clusters and a large parent
intensity with small ones produce point patterns whose second-order behaviour is
near-indistinguishable. Minimum contrast is fitted by numerical optimization of
a discrepancy that is close to flat along this trade-off, so the estimates it
returns in this regime are sensitive to the starting point and to the choice of
$r_{max}$. We deliberately retain this regime in our design distribution (see
Section \ref{sec:simulator}) rather than restricting to well-separated clusters,
and report error as a surface over the design grid so that the degradation is
visible rather than averaged away.

% --- (iii) the topological turn ---------------------------------------------
These limitations motivate the use of persistence-based features which provide a
multi-scale, process-agnostic descriptor of the point cloud that does not require
a closed-form or a choice of summary statistic. \todo{Not strictly true. Rephrase}

% --- (iv) contributions ------------------------------------------------------
\subsection{Contributions}
\label{ssec:contributions}

\emph{Notes to expand into prose; subsection not final}
\begin{itemize}
    \item \textbf{Estimator.} One architecture and training recipe -- persistence
    images of the $H_0$ diagram under a single $DTM_5$ filtration, a small CNN
    encoder, and concatenation with $\log N(\pc)$ -- designed
    to cover Thomas and nested Thomas; on Thomas it is
    competitive with (not decisively better than) a specialized
    classical-summary neural baseline, but unlike minimum contrast it
    transfers unchanged to nested Thomas, where no closed-form summary
    statistic exists to fit against.
    \item \textbf{Calibration methodology.} Per-diagram coverage-quantile
    calibration for summarizing across a heterogeneous population of
    diagrams, with an explicit train/test leakage protocol
    (Section~\ref{ssec:calibration}) and a sweep over $q$, $\alpha$,
    $\sigma_{\text{pixels}}$, $G$ showing which choices actually matter
    (Appendix~\ref{app:calib-pi}).
    \item \textbf{Validated simulator and evaluation protocol.} Multi-family
    simulator under a common
    reparameterization, validated against an exact finite-window second-order
    identity (Section~\ref{sec:simulator}); an evaluation protocol anchored
    between an upper, chance-level reference (label shuffling,
    Table~\ref{tab:label-shuffle}) and a lower, irreducible-error-floor
    reference ($\mathcal{L}_{\text{floor}}$, Section~\ref{sec:problem} --
    currently unresolved, see Section~\ref{ssec:exp-performance}).
    \item \textbf{What the estimator turns out not to need.} Multi-scale
    fusion, the $H_1$ channel, coordinate channels, and three alternative
    vectorizations are each reported as branches off the main line
    (Section~\ref{ssec:branches}); none of them earns a place in the model,
    and the single configuration where the simplification costs something --
    fused scales on nested Thomas -- is reported with them rather than
    absorbed.
\end{itemize}


% #############################################################################
% 2. RELATED WORK
% #############################################################################
% NOTE -- PURPOSE: concede the architecture to [A]; concede topological features
% for parameter estimation to [B]; claim the comparative study, the calibration
% methodology, and the multi-family scope.
% NOTE -- outstanding: three paragraphs -- (i) neural estimators for SPPs, [A]
% and its lineage; (ii) topological summaries as features for inference, [B]
% plus the vectorization literature; (iii) the gap, in one or two sentences.
% #############################################################################

\section{Related Work}
\label{sec:related}

REF [Vihrs 2022] first proposed the use of deep learning for parameter estimation. They implement
an encoder/inference architecture similar to PointNet, but use the $L$-function as the learned
feature and therefore suffer from the same limirations as classical methods (i.e. reliance
on a closed-form summary statistic). REF employ a similar architecture as REF, but actually use
topological features. Their model is deigned to estimate cosmological parameters rather than the
parameters of the underlying point process. Furhtermore, they do not systematically compare
various topological summaries, nor do they address the calibration of the summaries they use.

The vectorization of persistence diagrams is an active area of research arising from the issue
that persistence diagrams -- not being a Hilbert space -- cannot be direclty used as inputs for
machine learning models. To mitigate this, REF [Adams 2017] propose persistence images, which vectorise persistence
diagrams by applying a Gaussian kernel to each point in the diagram and intergating over a grid.
REF propose Betti curves, which vectorize persistence diagrams by counting the number of topological features
that are alive at each scale.


% #############################################################################
% 3. BACKGROUND
% #############################################################################
% NOTE -- PURPOSE: the SPP and persistent-homology material the rest of the
% paper assumes, and nothing else.
% #############################################################################

\section{Background}
\label{sec:background}

We provide the background in SPPs and persistent homology required for this
discussion. Readers with a working understanding of both may skip to
Section \ref{sec:problem}.

% -----------------------------------------------------------------------------
\subsection{Spatial point processes}
\label{ssec:spp}
% -----------------------------------------------------------------------------

    \begin{definition}[Intensity function]
    \label{def:intensity}
        A $d$-dimensional SPP $X$ on a subset $W \subseteq \R^d$ is equipped
        with an intensity function $\rho(u): \R^d \rightarrow \R^+$ such that
            \[
            \E[N(W)] = \int_W \rho(u) du,
            \]
        where $N(W) = \mid W \cap X \mid$ denotes the number of points of $X$
        in $W$. When $\rho(u) = \lambda \in \R$, $X$ is said to be
        \emph{homogeneous}.
    \end{definition}

    \begin{example}
        For a homogeneous Poisson process with intensity $\lambda$, we have
        $\rho(u) = \frac{1}{\lambda}$, i.e. points are scattered uniformly on
        the space.
    \end{example}

    \begin{definition}[Neyman-Scott process]
    \label{def:ns}
        Let $P$ be a homogeneous Poisson \emph{parent} process with intensity
        $\kpar$. For a Neyman-Scott process with Poisson-distributed offspring
        we have:
            \[
            \rho_{NS}(u) = \mo \sum_{p \in P} k(u - p),
            \]
        where $k$ is a kernel that determines how the \emph{offspring} points
        are distributed around the parents.
    \end{definition}

    \begin{note}
        Neyman-Scott processes always have the parent and offspring intensities
        as parameters. Additional parameters are introduced by the choice of
        kernel.
    \end{note}

    \begin{definition}[Homogeneous Thomas process]
    \label{def:thomas}
        A homogeneous Thomas process is a Neyman-Scott process with an isotropic
        Gaussian kernel $\mathcal{N}(0, \sigma^2I)$
            \[
            \rho_{Thomas}(u) = \mo \sum_{p \in P} k(0, \sigma^s),
            \]
        This process has parameters $\theta_{Thomas} = (\kappa, \mu, \sigma)$
        corresponding to the parent intensity, offspring intensity, and cluster
        scale, respectively.
    \end{definition}

% -----------------------------------------------------------------------------
\subsection{Persistent homology}
\label{ssec:persistence}
% -----------------------------------------------------------------------------

Radially averaged statistics -- such as the $g$ and $K$ functions -- collapse
a point cloud to a real-valued curve indexed by scale $r$. Features based on such
statistics are constructed by evaluating over a chosen range $[r_{min}, r_{max}]$
and subsampling at $n$ points along the interval to obtain a feature vector in $\R^n$.
This approach is limited because it discards multi-scale connectivity information and
void structure, i.e. \emph{how} the density is distributed across the space.
Persistent homology retains this information by tracking a family of simplicial
complexes on the point cloud as the scale parameter grows, instead of a single scalar
summary. Under a filtration such as Vietoris-Rips, the $H_0$ persistence summary records
the order and scales at which clusters merge. The $H_1$ counterpart marks voids and gaps
in the pattern as they persist across scales, allowing one to capture the clustering-induced
emptiness or repulsion that a process may exhibit, which canonical statistics cannot see as clearly
due to radial averaging.
For each $h \in \N$, the $h$-dimensional persistence summary requires no closed form expression
and no assumption about which process family generated it.

Throughout this section, let $X$ be a point process on $R^d$ and $W \subseteq \R^d$ a bounded window with
$|X \cap W| = n$ . Enumerate the points of $X \cap W$ to construct a point cloud
$\pc = \{x_1, ..., x_n\} \subset \R^d$. Let $\Sigma(\pc)$ denote the \emph{full
simplex} on $\pc$, i.e. the collection of all non-empty subsets of $\pc$.

    \begin{definition}[Simplex]
    \label{def:simplex}
        A \emph{simplex} of $\pc$ is a non-empty finite subset
        $\sigma \subseteq \pc$. Its faces are the non-empty subsets
        $\tau \subseteq \sigma$.
    \end{definition}

    \begin{definition}[Simplicial complex]
    \label{def:simpcomplex}
        A simplicial complex $\simpcom$ on $\pc$ is a collection of simplices
        closed under taking faces: $\sigma \in \simpcom, \tau \subseteq \sigma
        \implies \tau \in \simpcom$.
    \end{definition}

    \begin{definition}
    \label{def:filtfunc}
        A filtration function on $\pc$ is a monotone map
        $f: \Sigma(\pc) \rightarrow \R^+$ such that
            \[
            f(\tau) \leq f(\sigma) \quad \forall \quad \tau \subseteq \sigma.
            \]
    \end{definition}

    \begin{example}
        Vietoris-Rips filtration function:
            \[
            f_{VR}(\{x_i\}) = 0, \qquad f_{VR}(\{x_i, x_j\}) = \|x_i - x_j\|,
            \]
        extended to higher simplices by $f_{VR}(\sigma) = \max_{\{x_i,x_j\}
        \subseteq \sigma} f_{VR}(\{x_i, x_j\})$, i.e. the diameter of $\sigma$.
    \end{example}

    \begin{example}
    \label{ex:dtm}
        For $u \in \R^d$, let $N_k(u)$ denote its $k$ nearest neighbours in
        $\pc$. The $\mathrm{DTM}_k$ filtration function is
            \[
            f_{DTM}(\sigma; k) = \max_{u \in \sigma}\{(\frac{1}{k}
            \sum_{v \in N_k(u)} d(u, v)^2)^{\frac{1}{2}}\}
            \]
        again extended to higher simplices by the max-over-edges rule of the VR
        example. \citep{anai2020dtm}.
    \end{example}

    \begin{claim}
    \label{claim:dtm-monotone}
        $f_{DTM_\kappa}$ is monotone, and is hence a valid filtration function
        \citep[Prop.\ 3.5]{anai2020dtm}.
    \end{claim}

    \begin{definition}
    \label{def:filtration}
        Given a filtration function $f$ on $\pc$ and $r \geq 0$, we construct
            \[
            \mathcal{F}_r = \{\sigma \in \Sigma(\pc) : f(\sigma)
            \leq r\}.
            \]
        Monotonicity of $f$ ensures each $\mathcal{F}_r$ is a simplicial complex
        and that $\mathcal{F}_s \subseteq \mathcal{F}_t$ for $0 \leq s \leq t$.
        The family $(\mathcal{F}_r)_{r \geq 0}$ is the \emph{filtration} induced
        by $f$. We write $VR_r(\pc)$ for $\mathcal{F}_r$ under $f_{VR}$.
    \end{definition}

    \begin{definition}[Betti number]
    \label{def:betti}
        Fix $h \in \N$ and a simplicial complex $\Sigma$. Taking homology
        with coefficients in a field (we use $\Z_2$), $H_h(\simpcom)$ is a
        vector space of dimension $\beta_h = \dim H_h(\simpcom)$, the $h$-th
        \emph{Betti number}. $\beta_0, \beta_1, \beta_2, \dots$ count connected
        components, loops, and cavities, respectively.
    \end{definition}

    \begin{definition}[Persistence pair \& lifetime]
    \label{def:pers-pair}
        For a filtration $(\filt_r)_{r \geq 0}$, a $h$-dimensional homology
        class is \emph{born} at the smallest scale $b \geq 0$ at which it
        appears in $H_h(\simpcom_b)$, and \emph{dies} at the scale $d > b$ at
        which it merges into an older class or becomes a boundary. The pair
        $(b,d)$ is a \emph{persistence pair}, and $d - b$ its \emph{lifetime}.
    \end{definition}

    \begin{example}
        Consider $H_0$. Every $0$-dimensional class is born at $r=0$ and dies
        when its two components merge, i.e. at the filtration value of the
        connecting edge.
    \end{example}

    \begin{definition}[Persistence diagram]
    \label{def:pd}
        A persistence diagram $D_h = \{(b_i, d_i)\}_i$ is the multiset of
        dimension-$h$ persistence pairs, plotted with birth on the $x$-axis and
        death on the $y$-axis.
    \end{definition}

    \begin{remark}
        Every point of $D_h$ lies above the diagonal $b=d$ because it must die
        after it is born. Distance from the diagonal is lifetime, so short-lived
        points near the diagonal are often treated as noise.
    \end{remark}

    \begin{definition}[Tent function]
    \label{def:tent}
        For a persistence pair $(b_i, d_i)$, the tent function is
            \[
            \Lambda_i(t) = \max\bigl(0,\, \min(t - b_i,\, d_i - t)\bigr),
            \quad t \in \R,
            \]
        a triangular bump supported on $[b_i, d_i]$, peaking at height
        $(d_i-b_i)/2$ when $t = (b_i+d_i)/2$. Persistence landscapes and
        silhouettes are order statistics and weighted averages of this family;
        persistence images replace the tent with a smooth Gaussian bump over
        the same birth-persistence coordinates.
    \end{definition}

    \begin{definition}[Persistence image]
    \label{def:per_image}
        The persistence image of $D_h$ is constructed by:
            \begin{enumerate}
                \item Mapping to birth-persistence coordinates:
                      $(b,d) \mapsto (b, d-b)$.
                \item Weighting each point by $w(d-b)$.
                \item Forming the persistence surface
                      $\rho(z) = \sum_i w(d_i - b_i)\,
                      \phi(z; (b_i, d_i - b_i))$, a weighted sum of isotropic
                      Gaussian kernels with bandwidth $\sigma_{im}$ centered at
                      the transformed points.
                \item Integrating $\rho$ over a $64 \times 64$ pixel grid,
                      giving a single-channel image $PI \in \R^{64 \times 64}$.
            \end{enumerate}
    \end{definition}


% #############################################################################
% 4. PROBLEM SETUP
% #############################################################################
% NOTE -- PURPOSE: state the single-realization estimation problem, the training
% loss, and the irreducible error floor the results are read against.
% #############################################################################

\section{Problem Setup}
\label{sec:problem}

Let $X$ be a Neyman-Scott process on $\R^d$ with unknown parameter vector
$\theta \in \R^{n_{params}}$. For example, the homogeneous Thomas process has 
$\theta = (\kappa, \mu, \sigma)$. We observe a single realization of $X$ on a 
bounded window $W \subset \R^d$, giving a point cloud
$\pc = \{x_1, \dots, x_n\} \subset \R^d$ with $n = |X \cap W|$. The objective is to 
produce an estimate $\hat\theta = \hat\theta(\pc)$ of $\theta$ from $\pc$ alone.
We stress \emph{single realization}: $\theta$ is not recovered from repeated draws of $X$, nor 
is $\pc$ pooled or averaged across realizations. This is the setting a practitioner
faces: one observed pattern and one chance to estimate the generating parameters.
It is also what makes the problem hard in a way no estimator can fix: a 
fixed $\theta$ already induces a distribution over point clouds, so even the true 
parameter value gives rise to a spread of plausible $\pc$, and it is this irreducible 
spread —- not any weakness of a particular estimator -— that sets the error floor 
characterized below.

The model is trained on a set of synthetic clouds
$\{\pc_i, \theta_i\}_{i=1}^{N_{train}}$ by backpropagating the RMS loss
    \[
    \mathcal{L}(\hat\theta_i, \theta_i) = \sqrt{\sum_{j=1}^{n_{params}}(\hat\theta^j_i - \theta^j_i)^2},
    \]
where $\theta_i = (\theta_i^1, \dots, \theta_i^{n_{params}})$. The scale of each
parameter varies \todo{For example...}; therefore -- if not normalized -- the 
large-scale parameters will
dominate this loss. To adress this, we train in log-normalized parameter space,
    \[
    \{\pc_i, \theta_i\}_{i=1}^{N_{train}} \rightarrow
    \left\{\pc_i, \frac{\log(\theta_i) - \mu_{\log\theta}}
    {\sigma_{\log\theta}}\right\}_{i=1}^{N_{train}},
    \]
where -- abusing notation -- $\mu_{\log\theta}, \sigma_{\log\theta}$ denote the
mean and standard deviation of the logged parameters over the training set.

We draw one realization per parameter vector. The dispersion of $\hat\theta$
across realizations sharing a common $\theta$ gives the irreducible error floor
against which the model's performance should be read. We name this quantity
$\mathcal{L}_{\text{floor}}$: fix $\theta$, draw $R$ independent realizations
$\pc^{(1)}, \dots, \pc^{(R)}$, fit $\hat\theta^{(r)}$ to each with a fast
classical estimator (mincontrast by default), and compute $\mathcal{L}$
(Section~\ref{ssec:eval-protocol}) between $\{\hat\theta^{(r)}\}_{r=1}^R$ and
the shared true $\theta$; averaging over a grid of $\theta$ spanning the
design distribution gives $\mathcal{L}_{\text{floor}}$ itself. No estimator
that sees only a single realization can average below
$\mathcal{L}_{\text{floor}}$ over the same design distribution, since
$\mathcal{L}_{\text{floor}}$ is realized by an estimator that already knows
which $\theta$ generated the cloud and is degraded only by which realization
of it was drawn. This is the lower reference point for every result in
Section~\ref{sec:experiments}; the upper one, $\mathcal{L}\approx1$, is
defined in Section~\ref{ssec:eval-protocol}. Section~\ref{ssec:exp-performance}
reports where the empirical estimate of $\mathcal{L}_{\text{floor}}$
currently stands.

% #############################################################################
% 5. SIMULATOR AND EXPERIMENTAL DESIGN
% #############################################################################
% NOTE -- PURPOSE: how the training data is generated, and the evidence that it
% is generated correctly. The splits and the loss now live in Section 6, since
% they define the quantities the method section relies on.
% #############################################################################

\section{Simulator and Experimental Design}
\label{sec:simulator}

Training data consists of pairs $(\pc_i, \theta_i)$ in which $\theta_i$ is drawn
from a design distribution over parameters space and each $\pc_i$ is a
single realization of the corresponding process. We describe the design
distribution, sampling algorithm, and verification procedure used to generate
the synthetic data.
% -----------------------------------------------------------------------------
\subsection{Design distribution}
\label{ssec:design}
% -----------------------------------------------------------------------------

The naive approach of sampling $(\kappa, \mu, \dots)$ uniformly and
independently from intervals is inadequate. Point processes can
typically be separated into qualitative regimes of interest, which occupy curved
subsets of the parameter grid. Therefore uniform sampling would allocate most of
its mass to configurations that are either near-Poisson or degenerate.
We therefore sample in reparameterized coordinates that carry the same meaning
for every family considered: \todo{Condense, make general}
    \begin{align*}
        \lambda &= \kappa\mu && \text{expected offspring intensity,} \\
        \kappa & && \text{parent intensity,} \\
        \nu &= r\sqrt{\kappa} && \text{dimensionless overlap index,}
    \end{align*}
where $r$ denotes the cluster scale of the family in question: $r = \sigma$ for
Thomas and $r = \sqrt{\sigma_1^2 + \sigma_2^2}$
for nested Thomas. The overlap index $\nu$ measures cluster radius against mean
inter-parent spacing, so $\nu \ll 1$ gives well-separated clusters and
$\nu \gtrsim 1$ gives clusters that merge (the realization approaches complete
spatial randomness). Nested Thomas carries an additional coordinate, the scale
ratio $\rho = \sigma_2/\sigma_1 \in (0,1)$, which separates configurations in
which both levels of structure are resolvable from those in which the process
collapses to a single scale. \todo{Make general. more coords?}

We draw $(\log\kappa, \log\lambda, \log\nu)$ uniformly over a grid and invert to
recover $\mu = \lambda/\kappa$ and $r = \nu/\sqrt{\kappa}$ for the training labels. 
Log-uniform sampling
is used because the parameters span several orders of magnitude, and it matches
the log-normalization applied to the ground-truth targets (see
Section \ref{sec:problem}) so that training targets are approximately uniform in
each coordinate. Draws are taken from a scrambled Sobol sequence. \todo{what is sobol?}

The following constraints are applied to the design grid:
    \begin{enumerate}
        \item \textbf{Point budget.} The cost of computing persistence diagrams
        grows rapidly with the number of points, so we impose
        $\mathbb{E}[N] = \lambda \in [150, 800]$, fixing the range of
        $\log\lambda$.
        \item \textbf{Clusters per window.} We require $\kappa|W| = \kappa \geq 15$
        so that a single realization is informative about $\kappa$, and
        $\kappa|W| = \kappa \leq 120$ so that individual clusters remain
        resolvable.
        \item \textbf{Overlap.} We impose $\nu \in [0.10, 0.90]$. The upper
        limit deliberately admits the near-Poisson regime so that the model is
        still exposed to configurations in which the parameters are weakly
        identifiable.
    \end{enumerate}
Parameter draws violating these constraint are rejected and redrawn before any
cloud is generated. This filter acts on the design and not on realizations,
therefore it does not distort the conditional law of $\pc$ given $\theta$. The
realized acceptance rate is $24.8\%$.

% -----------------------------------------------------------------------------
\subsection{Sampling algorithm}
\label{ssec:sampling}
% -----------------------------------------------------------------------------

The instances of the Neyman-Scott construction may be generated with the following
general framework. The components that change from instance-to-instance is the
displacement kernel $k$ and the offspring count distribution. For a window
$W \subset \R^d$:
    \begin{enumerate}
        \item Expand $W$ by an edge buffer $b$ to produce $W_{exp}$, so that
        clusters whose parents lie outside $W$ are not truncated.
        \item Sample the parent process on $W_{exp}$:
        $n \sim \mathrm{Poisson}(\kappa|W_{exp}|)$, with parents placed uniformly
        on $W_{exp}$.
        \item For each parent $p$ independently, draw an offspring count
        $n_p \sim \mathrm{Poisson}(\mu)$.
        \item For each offspring independently, draw a displacement
        $\epsilon \sim k$ and place the point at $p + \epsilon$.
        \item Discard the parents and retain the offspring lying in $W$.
    \end{enumerate}
For the Thomas process $k = \mathcal{N}(0, \simpcom^2 I_d)$; for the Mat\'ern
cluster process $k$ is uniform on the ball of radius $R$; for nested Thomas the
parent process at step 2 is itself a Thomas process rather than a homogeneous
Poisson process.

The buffer is set by a single rule applied to whichever kernel is in use: $b$ is
the radius containing all but $\varepsilon$ of the mass of $k$. For a Gaussian
kernel this gives $b = 4\sigma$, at which the per-axis tail mass is
$2\Phi(-4) \approx 6.3 \times 10^{-5}$; for the compactly supported Mat\'ern
cluster kernel it gives $b = R$ exactly, with no tail margin required; for nested
Thomas the buffers of the two levels add, giving $b = 4\sigma_1 + 4\sigma_2$,
which over-covers the tight bound $4\sqrt{\sigma_1^2 + \sigma_2^2}$ and is
therefore conservative. Points are never clipped to the boundary; the only
boundary operation is the containment filter at step 5.
\todo{Explain buffer better}

% -----------------------------------------------------------------------------
\subsection{Validation}
\label{ssec:verification}
% -----------------------------------------------------------------------------

Throughout this section, let $X$ denote a point process on $\R^d$. We write
$u, v \in \R^d$ for points (precise locations) and $du$, $dv$ for infinitesimal
regions of volume $|du|$, $|dv|$ centered at $u$ and $v$, respectively.

    \begin{definition}[Second-order product density]
    \label{def:rho2}
        The second-order product density
        $\rho^{(2)}: \R^d \times \R^d \rightarrow \R^+$
        is defined such that $\rho^{(2)}(u, v)\,du\,dv$ is the probability of
        jointly observing a point of $X$ in each of the infinitesimal regions
        $du$ and $dv$.
    \end{definition}

    \begin{definition}[Pair correlation function]
    \label{def:pcf}
        The pair correlation function $g(u, v): R^d \times R^d \rightarrow \R^+$
        is defined
        \[
        g(u, v) = \frac{\rho^{(2)}(u, v)}{\rho(u)\rho(v)}.
        \]
    It is clear that for any Poisson process we have $g \equiv 1$. Intuitively,
    $g$ can suggest whether SPPs are attractive ($g > 1$), or repulsive
    ($g < 1$).
    \end{definition}

    \begin{definition}[Stationarity \& isotropy]
    \label{def:stat-iso}
        A point process $X$ is said to be:
            \begin{itemize}
                \item \emph{Stationary} if $X \sim X + s$ for every
                $s \in \R^d$ where $X + s = \{x + s: x \in X\}$.
                \item \emph{Isotropic} if $X \sim RX$ for every rotation $R$
                about the origin.
            \end{itemize}
    \end{definition}

    \begin{note}
        Stationarity implies constant intensity $\rho(u) = \lambda$, but the
        converse is not necessarily the case. If $X$ is stationary and isotropic,
        $g(u, v)$ depends on $u$ and $v$ only through their inter-point distance
        $r = \|u - v\|$. Abusing notation, we may write $g(r)$ for this common
        value. We now define the $K$ function.
    \end{note}

    \begin{definition}[$K$-function]
    \label{def:kfunc}
        For stationary and isotropic $X$, the $K$-function
        $K(r): \R^+ \rightarrow R^+$ is defined
            \[
            K(r) = \int_{b_r(0)} g(\|h\|) dh
            \]
        Intuitively, $\rho(u)K(r)$ is the expected number of points we expect to
        see within a distance $r$ of a typical point in $X$.
    \end{definition}

% --- validation proper -------------------------------------------------------
The validity of the proposed pipeline is dependent on a diligent cloud-generation
process, which we validate against the established closed-form second-order theory 
(where possible).
On $W = [0,1]^2$ we compare $3000$ realizations at each of three parameter
triples spanning the design, chosen to span the overlap index: a tight-cluster
regime ($\nu = 0.07$), a moderate regime ($\nu = 0.22$), and a near-Poisson
regime ($\nu = 1.00$). \todo{Explain better}

For the first moment, $\mathbb{E}[N(W)] = \kappa\mu|W|$. For the second, the
canonical identity $\mathrm{Var}[N(W)] = \kappa\mu(1+\mu)|W|$ holds only in the
limit of a window large relative to $\sigma$, and is badly biased in precisely
the near-Poisson regime we wish to stress. Applying the law of total variance to
$N(W) = \sum_p N_p(W)$ with
$N_p(W) \mid p \sim \mathrm{Poisson}(\mu q(p))$, $q(p) = \int_W k(x-p)\,dx$, and
evaluating the shot-noise term by Campbell's formula gives the exact
finite-window result \todo{Explain better}
    \begin{equation}
    \label{eq:varN}
    \mathrm{Var}[N(W)] = \kappa\mu + \frac{\kappa\mu^2}{4\pi\sigma^2}J(\sigma)^2,
    \qquad
    J(\sigma) = \int_{-1}^{1}(1-|t|)\,e^{-t^2/4\sigma^2}\,dt,
    \end{equation}
for $W = [0,1]^2$, where the reduction to a one-dimensional integral follows from
the triangular density of the difference of two independent
$\mathrm{Uniform}(0,1)$ variables. As $\sigma \to 0$,
$J(\sigma) \to 2\sigma\sqrt{\pi}$ and \eqref{eq:varN} collapses to
$\kappa\mu(1+\mu)$, confirming the two agree in the appropriate limit; it is
\eqref{eq:varN}, and not the canonical identity, that we use as the comparison
target.

\begin{table}[tbp]
\centering
\begin{tabular}{lccccc}
\toprule
 & & \multicolumn{2}{c}{$\mathbb{E}[N(W)]$} & \multicolumn{2}{c}{$\mathrm{Var}[N(W)]$} \\
\cmidrule(lr){3-4}\cmidrule(lr){5-6}
Case & $\nu$ & theory & empirical & theory & empirical \\
\midrule
A & 0.07 & 200.0 & $200.4 \pm 0.6$ & 982.0 & $972.5 \pm 25.1$ \\
B & 0.22 & 100.0 & $100.2 \pm 0.4$ & 545.2 & $544.8 \pm 14.1$ \\
C & 1.00 & 20.0  & $20.0 \pm 0.1$  & 43.6  & $46.0 \pm 1.2$ \\
\bottomrule
\end{tabular}
\caption{Simulator validation against \eqref{eq:varN}. All deviations are within
$2$ standard errors.}
\label{tab:sim-validation}
\end{table}

We additionally verify that the empirical pair correlation function and Ripley's
$K$ match
    \[
    g(r) = 1 + \frac{e^{-r^2/4\sigma^2}}{4\pi\kappa\sigma^2},
    \qquad
    K(r) = \pi r^2 + \frac{1 - e^{-r^2/4\sigma^2}}{\kappa},
    \]
that the marginal offspring intensity on $W$ shows no boundary deficit
(edge-to-interior density ratio $1.003$--$1.048$ across the three cases, in all
cases consistent with unity and never in the direction an insufficient buffer
would produce), and that the displacement kernel is isotropic with per-axis
standard deviation matching $\sigma$. Estimates of $g$ normalized per realization
carry a negative bias that grows with $\mathbb{E}[N^2]$ and is therefore larger
for clustered patterns than for a matched-intensity Poisson control; we account
for this when reading residual deviations rather than attributing them to the
sampler.


% #############################################################################
% 6. METHOD: PARAMNET
% #############################################################################
% NOTE -- PURPOSE: describe the model that is actually being proposed -- one
% filtration, one scale, one homology dimension, one encoder -- together with
% the protocol its numbers are reported under. Everything tried and dropped is
% in Section 7.3, not here.
% #############################################################################
\section{Method: ParamNet}
\label{sec:method}

Let $\pc$ denote a point cloud, $v$ its vectorized persistence summary, and
$e \in \R^{E + 1}$ the resulting feature vector. ParamNet follows the pipeline \todo{define E}
    \begin{equation}
    \label{eq:pipeline}
    \pc \xrightarrow{\;\text{Vectorization}\;} v
    \xrightarrow{\;\text{Encoder}\;} e
    \xrightarrow{\;\text{FCNN}\;} \hat\theta.
    \end{equation}

We note that the pipeline shape follows the neural-estimator patterns used in REF [Vihrs 2022]
and REF. Our contribution lies in the feature-extraction stage, the application to
a wider range of SPP families, and the careful evaluation of design choices.

The model is trained and evaluated on synthetic clouds generated according to the
methodology described in Section~\ref{sec:simulator}. The synthetic clouds are processed
together before training so that the feature-extraction step is performed only
once.

Feature extraction is the model's key component: it is where the cloud's
topological summary is vectorized and made usable for machine learning.
A substantial body of work has developed vectorizations of persistence
diagrams -- persistence images, Betti curves, landscapes, kernel methods --
that trade off stability, resolution, and information loss in different ways,
but there has been no firmalized discussion on which is preferable in this setting.

ParamNet uses the persistence image vectorization; the three alternatives we implemented,
and their relative efficacies are reported in Section~\ref{ssec:exp-summaries}.

The vectorized summary $v$ is passed through an encoder to obtain the
feature vector $e \in \R^{E + 1}$, which is then given to 
$NN_\phi: \R^{E + 1} \to \R^{n_{params}}$ to produce the estimate
$\hat\theta$. $NN_\phi$ is a small feed-forward network: 
two hidden layers of width 64 and 
32, each followed by a ReLU nonlinearity and dropout ($p=0.1$), then a final
linear layer mapping to the $n_{params}$ outputs.

Every experiment in this paper uses the same optimizer and schedule: AdamW ($\text{lr}=10^{-3}$,
weight decay $3\times10^{-4}$), batch size $32$, up to $500$ epochs with
early-stopping patience $50$ on the validation loss.

% -----------------------------------------------------------------------------
\subsection{Splits and evaluation protocol}
\label{ssec:eval-protocol}
% -----------------------------------------------------------------------------

Data are split by parameter vector so that no parameter
value appears in both training and evaluation. Beyond the random test split we
generate two further evaluation sets: a grid over the design grid, permitting
error to be reported as a surface over $(\log\nu, \log\lambda)$ rather than as a
single scalar; and a set drawn one increment outside each range, characterizing
degradation under extrapolation. \todo{Don't understand extra sets}

The classical baselines introduced in Section~\ref{sec:intro} are run on the same clouds,
under the same train/test/grid/extrapolation splits rather than on separate data
or numbers taken from the literature. Section~\ref{ssec:exp-baselines}
is therefore a comparison
against a shared  test set.

Throughout Section \ref{sec:experiments} we report error in the log-normalized
parameter space of Section \ref{sec:problem}. For a parameter vector with $n_{params}$
components and a test set of size $N_{test}$, we define the per-parameter loss
    \[
    \mathcal{L}^j = \frac{1}{N_{test}} \sum_{i=1}^{N_{test}}
    (\hat{\tilde\theta}^j_i - \tilde\theta^j_i)^2,
    \qquad
    \mathcal{L} = \frac{1}{n_{params}} \sum_{j=1}^{p} \mathcal{L}^j,
    \]
where $\tilde\theta$ denotes the log-normalized parameter vector and
$\hat{\tilde\theta}$ the model's estimate of it. Scalar figures quote
$\mathcal{L}$; where a single number obscures the behaviour of an individual
parameter we report the breakdown $\mathcal{L}^1, \dots, \mathcal{L}^{n_{params}}$
alongside it. \todo{Clean up}
 
Two remarks. Firstly, $\mathcal{L}$ is the training objective of Section
\ref{sec:problem} up to the averaging over parameters and the square root, so
the quantity being optimized and the quantity being reported differ only in
scale; we report $\mathcal{L}$ rather than the training loss because the
per-parameter decomposition is not available from the latter. Secondly, because
the standardization is the one used at training time, $\mathcal{L}^j$ is
measured in units of the training-set variance of $\log\theta^j$. A value of
$\mathcal{L}^j = 1$ therefore corresponds to an error of one standard deviation
of $\log\theta^j$ over the design distribution, which is the error of a model
that has learned nothing beyond the mean of the design distribution. This gives
the upper reference point for every figure reported in Section
\ref{sec:experiments}; the lower reference point is the irreducible error floor
of Section \ref{sec:problem}.

% -----------------------------------------------------------------------------
\subsection{Filtration and vectorization}
\label{ssec:filtration}
% -----------------------------------------------------------------------------

We describe the pipeline implemented for the
persistence image vectorization computed under the $DTM$ filtration, which is the primary
configuration carried through the rest of the paper. For an introduction to persistence
images see see Definition~\ref{def:per_image}.
    \begin{enumerate}
        \item Pick filtration function and construct filtrations
        $\{(\filt)_{r\geq 0}\}_i$.
        \item Compute persistence diagrams $\{D^{(h)}\}_i$.
        \item Calibrate diagrams to determine common birth/persistence axis
        ranges, shared across the dataset.
        \item Cectorize each diagram into a
        fixed-size raster using those ranges.
    \end{enumerate}
The key design choices that follow naturally from this pipeline are (1) the filtration function,
(2) the calibration procedure, and (3) the image resolution. In particular, the choice of filtration 
function plays a fundamental role in
what information is available to be extracted in the first place.

Instead of the canonical Rips or \v{C}ech filtrations, we use the $DTM_k$ filtration (see Example~\ref{ex:dtm}).
Rips is built from inter-point distances alone, making it sensitive to the
sparsest connection between two regions and does not register how densely those
regions are occupied. In this case $DTM_k$ is more suitable because it works with a point's average distance to its
$k$ nearest neighbours; consequently, $DTM_k$ provides a more direct measure of local density.
Since the point process parameters directly affect the local density of the processes,
we claim that this choice is natural. \todo{Calibration issues for Rips}

The parameter $k$ fixes the scale at which local density is read. Summaries
computed at several values of $k$ can be given to the encoder 
(with some additional architectural cnsiderations) to provide the model with fine and coarse
scales simultaneouly. We compute a single scale, $k=5$, noting that error grows
monotonically in $k$ on both Thomas and Nested Thomas. Fusing
$k \in \{5,10,15\}$ does not improve on $DTM_5$ on Thomas, though it does on
nested Thomas. Section~\ref{ssec:branch-multiscale} reports the relevant losses.

We compute the summary for $H_0$ only. $H_0$ records the order and scales at
which clusters merge -- the local density the parameters govern -- while $H_1$
contributes nothing measurable beyond it under this vectorization and lands at
chance level when it is used alone (Section~\ref{ssec:branch-homology}).

Three alternative vectorizations -- landscapes, silhouettes, Betti/Euler curves -- 
are defined and compared in Section~\ref{ssec:exp-summaries}.

% -----------------------------------------------------------------------------
\subsection{Calibration}
\label{ssec:calibration}
% -----------------------------------------------------------------------------

Consider the training clouds $\{\pc_i, \theta_i\}_{i=1}^{N_{train}}$
and fix the homology dimension $h \in \{0,1\}$ for some choice of filtration. We
obtain $N_{train}$ persistence diagrams $\{D^{(h)}_i\}_{i=1}^{N_{train}}$. Each
diagram will have a different distribution of points, some of which may have
extreme outliers with respect to the full set of diagrams.

To ensure that in the resulting persistence images each pixel bears the
same meaning, we calibrate the diagrams to have the same axis ranges. It
is not immediately clear that naive approach of calibrating to the maximum deaths 
and births (thus including all
points) is adequate in this context. This is because the resulting relative positioning of the
bulk of points: if a diagram in the training set has strong outliers, the remaining diagrams
will be rescaled in a way that places the majority of their points near the origin. 
Figure~\ref{fig:calibration_comparison}) demonstrates how this phenomenon might affect the vectorizations.

\begin{figure}[ht]
    \centering
    \includegraphics[width=0.46\textwidth]{calibration_comparison.pdf}
    \caption{Naive pooled min/max calibration (left) versus per-diagram
    coverage-quantile calibration, $q=0.95$, $\alpha=1.05$ (right), for the
    same $H_0$ diagram (DTM filtration, $k=5$, Thomas design distribution,
    Section~\ref{sec:simulator}). Top: the diagram against its calibrated
    birth-persistence box; the naive box is stretched by an outlier diagram
    elsewhere in the population, and the dashed square marks the region
    zoomed into below. Bottom: that region at each calibration's native
    pixel density -- a coarse $\sim\!5\times5$ block of mostly-fused pixels
    under the naive bound versus a finer $\sim\!11\times11$ grid resolving
    the secondary peak under the quantile bound, since the same
    $64\times64$ grid now covers a narrower range.}
    \label{fig:calibration_comparison}
\end{figure}

Our solution is to calibrate according to \emph{per-diagram coverage quantiles}.
For each training diagram $D_i^{(h)}$, define the three summary statistics
\begin{equation}
    \label{eq:calib-stats}
    b_i^{\min} = \min_{(b,d) \in D_i^{(h)}} b, \qquad
    b_i^{\max} = \max_{(b,d) \in D_i^{(h)}} b, \qquad
    p_i^{\max} = \max_{(b,d) \in D_i^{(h)}} (d - b),
\end{equation}
i.e. the smallest birth, largest birth, and largest persistence value
\emph{within a single diagram}. Given a coverage level $q \in (0,1]$ and a
padding factor $\alpha \geq 1$, we set the calibrated birth and persistence
ranges to be the $q$-quantiles of these per-diagram statistics, taken over the
training set:
\begin{equation}
    \label{eq:calib-ranges}
    \big[\,Q_{1-q}\big(\{b_i^{\min}\}_{i=1}^{N_{train}}\big),\
    \alpha \cdot Q_{q}\big(\{b_i^{\max}\}_{i=1}^{N_{train}}\big)\,\big]
    \qquad \text{and} \qquad
    \big[\,0,\ \alpha \cdot
    Q_{q}\big(\{p_i^{\max}\}_{i=1}^{N_{train}}\big)\,\big],
\end{equation}
where $Q_q(\cdot)$ denotes the empirical $q$-quantile. These two
ranges are shared by every diagram of homology dimension $h$, for every point
cloud in the dataset (train, validation, test, and adversarial alike), so that a
given pixel of the resulting persistence image always refers to the same region
of birth-persistence space.

This scheme is also compatible with the stability of persistence
images REF [Adams 2017], which are stated for a fixed bandwidth and a
fixed image domain: since $\sigma_{\text{pixels}}$ is specified in pixel units
(Definition~\ref{def:per_image}) rather than in absolute birth/persistence units,
a pixel remains a comparable, well-defined location across every image built from
the calibrated imager.

The calibrated ranges are fit \emph{exclusively} on $\{D_i^{(h)}\}_{i \in \text{train}}$, 
then frozen and applied unchanged to validation, test, and adversarial diagrams.
Fitting the ranges on the full, unsplit population of clouds leaks information 
about the held-out clouds into the training images. 

% As the train/validation/test partition is redrawn per random seed, calibration must
% likewise be refit per seed, on that seed's training rows only, and then applied
% frozen to the corresponding validation, test, and adversarial rows.

For homology dimension $H_0$, birth values are often identically zero (or
otherwise carry negligible spread) under the chosen filtration, in which case
$b_i^{\min} \approx b_i^{\max}$ for all $i$ and the calibrated birth range
collapses to a point. In these cases we fall back to a
one-dimensional vectorization since
the degenerate range would make an entire column of pixels convey no information.

The coverage-quantile rule is not specific to the persistence image: the
same $(q,\alpha)$ scheme is applied to the corresponding per-diagram statistics of
the 1-D grid underlying landscapes, silhouettes, and Betti/Euler curves 
(Appendix~\ref{app:calibration-sweeps}).

The four calibration parameters -- coverage $q$, padding $\alpha$, kernel width
$\sigma_{\text{pixels}}$, and raster resolution $G$ -- are swept one at a time
around default values, each arm varying a single parameter while holding
architecture and the remaining three fixed
(Appendix~\ref{app:calib-pi}). Only $G=128$ is unambiguously bad; $q$ is the one
parameter on which the two processes disagree (Thomas prefers $q=0.99$ and Nested Thomas
favors $q=0.95$), and we keep the shared
default $q=0.95$ rather than tune it per-process. \todo{Future work}

% -----------------------------------------------------------------------------
\subsection{Encoder and head}
\label{ssec:encoder-fusion}
% -----------------------------------------------------------------------------

Consider the training clouds $\{\pc_i, \theta_i\}_{i=1}^{N_{train}}$. For each
cloud we compute the $H_0$ persistence image of the $DTM_5$ diagram at
resolution $64 \times 64$. So for batch size $B$, the encoding stage is
comprised of the following pipeline:
    \begin{align*}
    (B, 1, 64, 64) \xrightarrow{\;\text{Encoder}\;} (B,E) \xrightarrow{\;\text{Concatenation}\;} (B,E+1)
    \end{align*}

In the encoder stage, the image is passed through a small
convolutional stack: three $3\times3$ convolution blocks with channel widths
$(32, 64, 128)$, each followed by a ReLU and $2\times2$ max-pooling with
spatial dropout ($p=0.2$) between blocks (except the last one). The resulting
feature map is reduced by adaptive average pooling to a single spatial location
and projected by a linear layer to the embedding dimension $E$ (we use $E=64$)
to produce a $(B, E)$ tensor.

The concatenation step appends $\log N(\pc)$, the logarithm of the number of points in
the cloud, to that embedding, giving a feature vector of dimension $E + 1$. The
resulting $(B, E + 1)$ tensor is then passed to the FCNN head to produce the
final parameter estimates.

With a single $DTM_k$ scale only one encoder is necessary, and there is no need for 
an architecture that fueses multiple scales. Additional complexitites arise when
several scales are computed. We must chose whether to use a shared encoder
(every scale is fed to an encoder with the same weights), or assign a seprate
encoder to each scale (thus increasing the number of encoder parameters threefold).
Once the encoder framework is determined, a suitable method for fusing the per-$k$
embeddings is necessary; for example simple concatenation, or more complex fusion
architecture. The training results obtained from various such combinations are 
reported in Section~\ref{ssec:exp-ablations}.


% #############################################################################
% 7. EXPERIMENTS
% #############################################################################
% NOTE -- PURPOSE: two jobs, in order. (1) what the model of Section 6 achieves
% and how to read that number, against the classical and neural baselines;
% (2) the branches -- everything implemented, run, and left out of the model.
% #############################################################################

\section{Experiments}
\label{sec:experiments}

We first report the estimation performance of the model described in
Section~\ref{sec:method}, together with the controls that make the numbers
interpretable, and compare it against existing classical and neural baselines on shared test
data. We then report the variations that were run and did not beat the current
iteration of ParamNet.

A joint sweep over calibration paramaters, architectures, and vectorizations
is computationally unfeasible. Instead, we adopt a greedy sequence fixed in advance.
Encoder and fusion mode are selected first
(Section~\ref{ssec:exp-ablations}), against the persistence image at the
default calibration ($q=0.95$, Section~\ref{ssec:calibration}). That
architecture is then held fixed while the calibration parameters
$(q,\alpha,\sigma_{\text{pixels}},G)$ are swept (Appendix~\ref{app:calib-pi}).
A confirmation pass re-runs the two architectures that survive the fusion-mode
sweep -- shared and independent encoders under concatenation -- at the
selected calibration (Table~\ref{tab:encoder-confirm}), to check that the
architecture choice was not an artifact of the default calibration under which it
happened to be made. Vectorization-specific parameters (landscape $K$,
silhouette $p$, etc.) are swept last, with calibration and architecture
both frozen at their selected values (Section~\ref{ssec:exp-summaries},
Appendix~\ref{app:calibration-sweeps}). \todo{Slit up, repetition from calib section}

Three conventions recur across every results table; we state them here to avoid
unecessary repetition in the captions. All reported losses are the
squared error in log-normalized space defined in Section~\ref{ssec:eval-protocol},
averaged over parameters. A dagger ($^\dagger$) marks an arm on
which at least one seed failed to converge and collapsed instead to the
$\mathcal{L}\approx0.9$--$1.0$ floor; in these cases the reported
mean is dominated by the failed seed(s) and should not be read as the arm's
typical performances. \todo{Further work}

% -----------------------------------------------------------------------------
\subsection{Estimation performance}
\label{ssec:exp-performance}
% -----------------------------------------------------------------------------

The model of Section~\ref{sec:method} is assembled from choices each of which
was measured one axis at a time against the multi-scale, $H_0{+}H_1$,
coordinate-channel configuration this work started from: dropping to $DTM_5$,
dropping $H_1$, and dropping the coordinate channels each land within one
seed-standard-deviation of the arm they were measured against on Thomas
(Section~\ref{ssec:branches}). No run configures all three at once, so the
figures quoted here and in Table~\ref{tab:classical-comparison} --
$\mathcal{L}=0.126\pm0.008$ on Thomas and $\mathcal{L}=0.196\pm0.006$ on nested
Thomas -- belong to that starting configuration rather than to the model as
specified, and the end-to-end confirmation is outstanding.
\expt{Train and evaluate the headline configuration exactly as Section~\ref{sec:method}
specifies it -- persistence image, $H_0$ only, $DTM_5$, shared encoder,
concatenation with $\log N$, no coordinate channels, $q=0.95$ -- on both design
distributions at $n=10$ seeds, so that the reported $\mathcal{L}$ is the model
the paper describes and not the configuration each simplification was measured
against.}

Two controls fix the scale on which those numbers should be read.

\paragraph{Sanity check: label shuffling.} The $^\dagger$ divergence in
Table~\ref{tab:calibration-sweep} is described as ``no better than predicting
the design-distribution mean,'' but that floor was never independently
measured. We do so directly: retraining the default PI model with training
labels randomly permuted (inputs unchanged), so the only learnable target is
each parameter's marginal distribution.
\begin{table}[tbp]
\centering
\small
\begin{tabular}{lcc}
\toprule
Arm & Thomas $\mathcal{L}$ & Nested-Thomas $\mathcal{L}$ \\
\midrule
Full model (correct labels) & $0.124 \pm 0.007$ & $0.196 \pm 0.006$ \\
Shuffled labels             & $1.011 \pm 0.024$ & $0.997 \pm 0.018$ \\
\bottomrule
\end{tabular}
\caption{Label-shuffle control, $n=5$ seeds, both design distributions. Same
architecture and calibration as the full model; only the training-label
assignment is permuted.}
\label{tab:label-shuffle}
\end{table}
Both datasets land at $\mathcal{L}\approx1.0$, confirming the $^\dagger$ rows
of Table~\ref{tab:calibration-sweep} are genuinely collapsing to
chance-level prediction and not to some other failure mode, and giving an
empirical floor against which every other result in this section can be
read.

\begin{table}[tbp]
\centering
\small
\begin{tabular}{lcc}
\toprule
Arm & Thomas $\mathcal{L}$ & Nested-Thomas $\mathcal{L}$ \\
\midrule
Full model (PI, $H_0{+}H_1$, DTM$\{5,10,15\}$) & $0.124 \pm 0.007$ & $0.196 \pm 0.006$ \\
$\log N$ removed & $0.305 \pm 0.392^\dagger$ & $0.200 \pm 0.006$ \\
Topology removed ($\log N$ only) & $0.891 \pm 0.017$ & $0.946 \pm 0.018$ \\
\bottomrule
\end{tabular}
\caption{Bracketing the topological contribution: the full model against a
model stripped of the $\log N(\pc)$ feature (Section~\ref{ssec:encoder-fusion})
and against a model with the topology removed entirely, keeping only $\log
N(\pc)$ fed straight to the FCNN head. Both design distributions, $n=5$ seeds
per arm. One of five Thomas seeds diverged to the design-distribution-mean
floor ($\mathcal{L}\approx1.0$, cf.\ Table~\ref{tab:calibration-sweep}); the
other four land at $0.123$--$0.133$, in line with the full model.}
\label{tab:logn-ablation}
\end{table}

Training on $\log N$ alone (i.e. no topological summary) collapses to the same
$\mathcal{L}\approx0.9$--$1.0$ floor as the label-shuffle control of
Table~\ref{tab:label-shuffle} on both datasets, so the point count carries
almost none of the recoverable signal on its own. Conversely, removing $\log N$ from the
full model costs little on nested-Thomas ($0.196\to0.200$) and is
merely noisy on Thomas rather than clearly harmful (four of five seeds are
indistinguishable from the full model; the fifth diverges). Together the two
arms show the persistence images -- not the point count -- are carrying the
signal, with $\log N$ a minor, possibly non-essential addition.

\paragraph{Performance relative to the irreducible error floor.}
Section~\ref{sec:problem} names $\mathcal{L}_{\text{floor}}$ -- the
dispersion of a fast classical estimator's $\hat\theta$ across many
realizations sharing one true $\theta$ -- as the lower reference point every
number in this section should ultimately be read against, complementing the
upper reference point $\mathcal{L}\approx1$ already carried as the bottom row
of Table~\ref{tab:classical-comparison}. $\mathcal{L}_{\text{floor}}$ itself
does not yet have a value to report.
\texttt{scripts/estimate\_error\_floor.py} was run on a dedicated
$50\,\theta \times 200$-realization Thomas dataset
(\texttt{data/error\_floor/thomas}, mincontrast estimator, 10 SLURM shards
covering all 50 $\theta$-groups): every shard completed without error
(1000/1000 clouds fitted per shard, confirmed in each shard's own log) but
every per-cloud loss reduced to \texttt{NaN} rather than a number, so all 50
groups -- and hence the merged \texttt{floor\_L} -- are \texttt{NaN}. This
looks like a bug in the script's own aggregation step rather than a property
of the estimator: the same \texttt{mincontrast} module, run on ordinary
(non-repeated-realization) clouds above, completes without crashing on all
5 seeds of Table~\ref{tab:classical-comparison} even where 2 of those 5 land
at elevated loss, so it is not silently producing \texttt{NaN} there. One
concrete lead: the shard JSONs' recorded \texttt{label\_names} include
\texttt{edge\_buffer} and \texttt{c1}, names that do not belong to Thomas's
own $(\kappa,\mu,\sigma)$ target set, suggesting the per-target loss array is
being indexed against the wrong label list for this dataset. Nested Thomas
cannot take the same route at all under the current codebase: it has no
closed-form $K$ or $g$ -- the same closed-form dependency
Section~\ref{sec:intro} raises against classical estimators generally -- so
\texttt{configs/runs/nested\_thomas/error\_floor.yaml} generates the
repeated-realization dataset but nothing yet turns it into a floor estimate
for that process. \expt{Fix the label-indexing bug in \texttt{estimate\_error\_floor.py}'s
aggregation and re-run for Thomas; nested Thomas needs either a derived closed
form or a different per-realization floor estimator (e.g.\ \texttt{vihrs}, which
needs no closed form) before the same route is even available.} Until
$\mathcal{L}_{\text{floor}}$ has a value, every loss in this section should
be read only against the upper, chance-level reference point.

\expt{Three evaluation sets defined in Section~\ref{ssec:eval-protocol} are not
yet reported and are each promised elsewhere in the paper: error as a surface
over $(\log\nu, \log\lambda)$ on the grid set -- expect degradation as
$\nu \to 1$, where the parameters are weakly identifiable, and if the model
degrades more gracefully than minimum contrast there, that IS the result;
degradation on the out-of-range set one increment outside each design range;
and the Mat\'ern cluster runs for the cross-family comparison -- whether one
model trained across families beats per-family models is an interesting question
in itself, and the common reparameterization of Section~\ref{ssec:design} is
what makes it askable.}

% -----------------------------------------------------------------------------
\subsection{Comparison against baselines}
\label{ssec:exp-baselines}
% -----------------------------------------------------------------------------

We compare ParamNet (PI, DTM$_{\{5,10,15\}}$, Section~\ref{ssec:branch-multiscale})
against the classical and neural baselines of Section~\ref{sec:intro}, run on
the identical clouds and train/test splits of Section~\ref{ssec:eval-protocol}
rather than against numbers taken from the literature: minimum contrast fit to
Ripley's $K$ and, separately, to the pair correlation function $g$
(Section~\ref{ssec:verification}), and \texttt{vihrs} -- our re-implementation
of the neural estimator conceded as prior work [A] in Section~\ref{sec:related},
Vihrs (2022): a 1-D CNN over the edge-corrected, centred Ripley $L(r)-r$ curve
plus $\log N(\pc)$, i.e.\ the pipeline of Section~\ref{sec:method} with the
persistence-image feature replaced by a single classical radial summary and no
persistence diagrams at all. Palm likelihood \citep{tanaka2008parameter} is
implemented (\texttt{cloudforger.baselines.palm}) but has not yet been run to
completion; \#\#\# marks it below.

\begin{table}[tbp]
\centering
\small
\begin{tabular}{lcc}
\toprule
Method & Thomas $\mathcal{L}$ & Nested-Thomas $\mathcal{L}$ \\
\midrule
ParamNet (ours)         & $0.126 \pm 0.008$ & $0.196 \pm 0.006$ \\
vihrs ($L(r)-r$ CNN)     & $0.115 \pm 0.007$ & $0.215 \pm 0.004$ \\
Minimum contrast, $K$   & $0.827 \pm 0.646^{a}$ & --- \\
Minimum contrast, $g$   & $6.167 \pm 0.331^{b}$ & --- \\
Palm likelihood         & \#\#\# & \#\#\# \\
\midrule
\emph{Chance-level reference (label shuffle, Table~\ref{tab:label-shuffle})} & \emph{$1.011 \pm 0.024$} & \emph{$0.997 \pm 0.018$} \\
\bottomrule
\end{tabular}
\caption{Estimator comparison, both design distributions, $n=10$ seeds
(ParamNet) or $n=5$ seeds (vihrs, minimum contrast) on the shared test split
of Section~\ref{ssec:eval-protocol}; \#\#\# = not yet run. Minimum contrast
has no nested-Thomas entry because the process has no closed-form $K$ or $g$
to fit against -- the same closed-form dependency Section~\ref{sec:intro}
raises against classical estimators in general, here making the comparison
itself impossible rather than merely inconvenient. The bottom row is the
upper (chance-level) reference point of Section~\ref{ssec:eval-protocol},
repeated from Table~\ref{tab:label-shuffle}, so every method above can be read
against both $0$ and chance. $^a$ Minimum contrast on $K$ is unstable across
seeds rather than uniformly bad: 3/5 seeds land at $0.379$--$0.399$, close to
vihrs, while 2/5 land far worse ($0.903$ and $2.056$, the latter driven by a
parent-intensity error of $3.24$) -- consistent with
Section~\ref{sec:intro}'s claim that minimum contrast is sensitive to the
optimizer's starting point specifically because the design distribution
retains the near-Poisson regime where $\kappa$ and $\mu$ are weakly
identifiable. $^b$ Minimum contrast on $g$ is, in contrast, consistently
worse than chance on $5/5$ seeds ($5.82$--$6.66$, tight relative to that
mean) -- a systematic bias rather than scattered non-convergence. Its own
module (\texttt{cloudforger.baselines.mincontrast\_g}) documents that its
contrast exponent $c=1$ is set by analogy with \citet{diggle2003spatial}'s
$K$-based rule of thumb rather than tuned, and that the estimator has not been
checked against ground-truth parameters the way the $K$-based fit implicitly
has been (Section~\ref{ssec:verification}); we report the number rather than
suppress it, but do not read it as a settled statement about contrast
estimation on $g$ in general.}
\label{tab:classical-comparison}
\end{table}

\texttt{vihrs} is not merely competitive with ParamNet on Thomas, it is
better -- and not only in aggregate. Its per-parameter losses
($\mathcal{L}^\kappa=0.194\pm0.012$, $\mathcal{L}^\mu=0.123\pm0.008$,
$\mathcal{L}^\sigma=0.029\pm0.002$) are each individually lower than
ParamNet's own ($0.204\pm0.010$, $0.132\pm0.009$, $0.041\pm0.006$;
Table~\ref{tab:filtration-comparison}'s DTM$_{\{5,10,15\}}$ row), so this is
not a scalar-averaging artifact: for this single-level process, a hand-built
radial summary plus a point count is not a straw-man baseline, and our
topological pipeline does not beat it. The picture is different on nested
Thomas, where ParamNet's advantage in Table~\ref{tab:classical-comparison}
is concentrated on the density/count targets --
$\mathcal{L}^{\text{parent\_intensity}}=0.286$ vs.\ vihrs's $0.314$,
$\mathcal{L}^{\text{meta\_offspring}}=0.420$ vs.\ $0.487$,
$\mathcal{L}^{\text{mean\_offspring}}=0.087$ vs.\ $0.111$ -- while
\texttt{vihrs} keeps a clear edge on both scale parameters,
$\mathcal{L}^{\text{meta\_cluster\_scale}}=0.134$ vs.\ ParamNet's $0.144$ and
$\mathcal{L}^{\text{cluster\_scale}}=0.029$ vs.\ $0.042$. The scale/count
split is consistent across both design distributions: $L(r)-r$ is built to
resolve a characteristic radius directly, while a density-based topological
feature is not, and it is only on the harder two-level process -- where
density and count are harder to read off a single radial curve -- that this
trade-off nets out in the topological pipeline's favour overall.

Both neural methods dominate minimum contrast in either variant, by roughly
an order of magnitude on Thomas and categorically on nested Thomas, where
minimum contrast cannot even be attempted. Minimum contrast on $g$ in
particular lands \emph{above} the chance-level floor of
Table~\ref{tab:classical-comparison}'s bottom row -- worse than predicting
the design-distribution mean -- which Section~\ref{sec:intro}'s critique of
classical methods did not anticipate needed stating so starkly; see the
table's own caption for the caveats around that number before over-reading
it.

\paragraph{Per-parameter breakdown.}
$\kappa,\mu,\sigma$ separately, for every method run on Thomas, are already
laid out in Table~\ref{tab:classical-comparison}'s discussion above and
Table~\ref{tab:filtration-comparison}'s DTM$_{\{5,10,15\}}$ row; the same
breakdown for the DTM$_k$ filtration sweep itself (rather than across
methods) is in Table~\ref{tab:filtration-comparison} directly. The
short version: $\mathcal{L}^\kappa$ (parent intensity) is the hardest target
for every method on both processes by a wide margin, $\mathcal{L}^\sigma$ /
$\mathcal{L}^{\text{cluster\_scale}}$ the easiest, and this ordering holds
regardless of which estimator is used -- so the difficulty ranking across
parameters looks like a property of the estimation problem itself (how much
each parameter's variation actually perturbs the observed point pattern)
rather than an artifact of any one feature set or architecture.

\begin{table}[tbp]
\centering
\small
\begin{tabular}{lccc}
\toprule
Method & $\mathcal{L}^\kappa$ (parent int.) & $\mathcal{L}^\mu$ (offspring int.) & $\mathcal{L}^\sigma$ (cluster scale) \\
\midrule
ParamNet (ours)        & $0.204 \pm 0.010$ & $0.132 \pm 0.009$ & $0.041 \pm 0.006$ \\
vihrs                  & $0.194 \pm 0.012$ & $0.123 \pm 0.008$ & $0.029 \pm 0.002$ \\
Minimum contrast, $K$  & $1.317 \pm 1.009$ & $0.912 \pm 0.726$ & $0.250 \pm 0.202$ \\
Minimum contrast, $g$  & $8.792 \pm 0.483$ & $6.286 \pm 0.355$ & $3.425 \pm 0.161$ \\
\bottomrule
\end{tabular}
\caption{Per-parameter breakdown of Table~\ref{tab:classical-comparison}'s
Thomas column ($n=10$ ParamNet, $n=5$ others). \texttt{vihrs} has the lowest
loss in every column; see the discussion above for why this does not carry
over to nested Thomas.}
\label{tab:classical-per-param}
\end{table}

% -----------------------------------------------------------------------------
\subsection{Branches}
\label{ssec:branches}
% -----------------------------------------------------------------------------

Each of the following was implemented and run, and none of it is part of the
model of Section~\ref{sec:method}. We report them in order of how close the
discarded option sits to that model.

\subsubsection{Multi-scale DTM fusion}
\label{ssec:branch-multiscale}

We also compute the persistence image at several $DTM$ scales at once and fuse
the resulting embeddings. On Thomas this does not improve on the single scale
$k=5$, so the headline model does not carry it; on nested Thomas it does
improve, and that exception is reported here rather than absorbed.

The $n_k$ scales are handled by folding the scale axis into the batch dimension
for a single forward pass -- one call over $B \cdot n_k$ images rather than
$n_k$ separate calls -- so the weights are tied across scales by construction
rather than by an explicit sharing mechanism. The resulting $(B, n_k, E)$ tensor
is fused by concatenating the $n_k$ embeddings along the feature axis before
$\log N(\pc)$ is appended, giving $(B, n_k \cdot E + 1)$.

\begin{table}[tbp]
\centering
\footnotesize
\begin{tabular}{lccccc}
\toprule
Filtration & $\mathcal{L}$ (mean $\pm$ sd) & $\mathcal{L}^\kappa$ & $\mathcal{L}^\mu$ & $\mathcal{L}^\sigma$ & Failed seeds \\
\midrule
Rips                          & ---$^a$            & ---               & ---               & ---               & $10/10^a$ \\
DTM$_5$                       & $0.125 \pm 0.006$  & $0.201 \pm 0.009$ & $0.134 \pm 0.007$ & $0.040 \pm 0.004$ & $0/10$ \\
DTM$_{10}$                    & $0.139 \pm 0.008$  & $0.221 \pm 0.016$ & $0.147 \pm 0.008$ & $0.049 \pm 0.003$ & $0/10$ \\
DTM$_{15}$                    & $0.146 \pm 0.006$  & $0.227 \pm 0.011$ & $0.155 \pm 0.008$ & $0.056 \pm 0.010$ & $0/10$ \\
DTM$_{\{5,10,15\}}$ (default) & $0.126 \pm 0.008$  & $0.204 \pm 0.010$ & $0.132 \pm 0.009$ & $0.041 \pm 0.006$ & $0/10$ \\
\bottomrule
\end{tabular}
\caption{Filtration comparison, Thomas process, persistence-image
vectorization at the default calibration ($q=0.95$, Section~\ref{ssec:calibration}),
$H_0{+}H_1$, shared-weight CoordConv architecture
(Section~\ref{ssec:encoder-fusion}), $n=10$ seeds per row (Rips: $10$
attempted). $\mathcal{L}^\kappa,\mathcal{L}^\mu,\mathcal{L}^\sigma$ are the
per-parameter breakdown (parent intensity, offspring intensity, cluster
scale respectively; Section~\ref{ssec:eval-protocol}). $^a$ Rips fails
identically on $10/10$ seeds \emph{before} training starts, not from a
convergence failure: under Rips every $H_0$ feature is born at $r=0$, so the
calibrated birth range collapses to a point and the calibration step raises
rather than silently degrading -- exactly the degenerate-axis failure mode
anticipated in Section~\ref{ssec:calibration}. No loss is produced for this
row by construction, and no $DTM$ row shows a failed seed.}
\label{tab:filtration-comparison}
\end{table}

\begin{table}[tbp]
\centering
\small
\begin{tabular}{lcc}
\toprule
Filtration & Thomas $\mathcal{L}$ & Nested-Thomas $\mathcal{L}$ \\
\midrule
Rips                          & ---$^a$           & ---$^a$ \\
DTM$_5$                       & $0.125 \pm 0.006$ & $0.205 \pm 0.003$ \\
DTM$_{10}$                    & $0.139 \pm 0.008$ & $0.237 \pm 0.007$ \\
DTM$_{15}$                    & $0.146 \pm 0.006$ & $0.260 \pm 0.016$ \\
DTM$_{\{5,10,15\}}$ (default) & $0.126 \pm 0.008$ & $0.196 \pm 0.006$ \\
\bottomrule
\end{tabular}
\caption{Table~\ref{tab:filtration-comparison}'s comparison repeated for both
design distributions, scalar $\mathcal{L}$ only -- nested Thomas's five
targets do not share Thomas's $\kappa,\mu,\sigma$ decomposition. $n=10$ seeds
per row, except nested-Thomas's DTM$_{\{5,10,15\}}$ row, which reuses the
$n=5$ \texttt{confirm\_shared\_concat} run already reported as the PI/native
entry of Table~\ref{tab:main-comparison} and the shared-encoder row of
Table~\ref{tab:encoder-confirm}, rather than a separately-seeded run.
$^a$ Rips fails identically on both design distributions, $10/10$ seeds each,
for the reason given in Table~\ref{tab:filtration-comparison}'s caption.}
\label{tab:filtration-comparison-nt}
\end{table}

Error grows monotonically with $k$ within a single scale, on both design
distributions: DTM$_5$ is the strongest individual scale and DTM$_{15}$ the
weakest (Thomas: $0.125\to0.146$; nested Thomas: $0.205\to0.260$, a
proportionally larger degradation), consistent with a finer neighbourhood
retaining more of the local-density signal the parameters govern while a
coarser one averages exactly that structure away. Fusing all three scales
tells two different stories across the two processes. On Thomas, the fused
arm ($0.126\pm0.008$) does not improve on the single best scale, DTM$_5$
($0.125\pm0.006$): the two are within one seed-standard-deviation of each
other, so fusion's advantage there is not raw accuracy but not having to
select $k$ in advance. On
nested Thomas the fused arm ($0.196\pm0.006$) clearly beats every single
scale tried, including DTM$_5$ ($0.205\pm0.003$) -- the two-level design
benefits from combining scales in a way the single-level one does not,
plausibly because the meta- and sub-cluster structure sit at genuinely
different characteristic radii that no single $k$ resolves together. Dropping
to the single scale is therefore free on Thomas and costs $0.009$ on nested
Thomas, and it is the one place where the model of Section~\ref{sec:method} is
a trade rather than a simplification.

Rips fails outright under this vectorization on both processes, for the
degenerate-birth reason given in Section~\ref{ssec:filtration}; every other
result in Section~\ref{ssec:exp-summaries} and
Appendix~\ref{app:calibration-sweeps} is accordingly reported only under
$DTM$.

\subsubsection{Homology dimensions}
\label{ssec:branch-homology}

We also compute the $H_1$ diagram alongside $H_0$. For the persistence image it
adds nothing measurable to $H_0$, and on its own it collapses to the
label-shuffle floor, so the headline model uses $H_0$ only.

\begin{table}[tbp]
\centering
\small
\begin{tabular}{llcc}
\toprule
Vectorization & Homology dims & Thomas $\mathcal{L}$ & Nested-Thomas $\mathcal{L}$ \\
\midrule
PI & $H_0$ only     & $0.126 \pm 0.006$ & $0.194 \pm 0.006$ \\
PI & $H_1$ only     & $0.895 \pm 0.019$ & $0.911 \pm 0.013$ \\
PI & $H_0{+}H_1$ (default) & $0.124 \pm 0.007$ & $0.196 \pm 0.006$ \\
\addlinespace
LS & $H_0$ only     & $0.135 \pm 0.006$ & $0.246 \pm 0.002$ \\
LS & $H_1$ only     & \#\#\# & \#\#\# \\
\addlinespace
BC & $H_0$ only     & $0.237 \pm 0.012$ & $0.294 \pm 0.014$ \\
BC & $H_1$ only     & $0.260 \pm 0.019$ & $0.510 \pm 0.009$ \\
\bottomrule
\end{tabular}
\caption{Homology-dimension ablation (C3), $n=5$ seeds per arm, both design
distributions, native encoder path at the shared $q=0.95$ calibration
(Section~\ref{ssec:calibration}) rather than each vectorization's own
appendix-preferred default. PI and BC are complete;
LS's $H_1$-only arm has not been run (\#\#\# = not yet run) -- LS's
$H_0{+}H_1$ default is the native row of Table~\ref{tab:main-comparison}
($0.132\pm0.008$ T / $0.250\pm0.005$ NT, 3 seeds), not repeated here since it
was not re-run at $n=5$.}
\label{tab:homology-ablation}
\end{table}
For the persistence image, $H_1$ alone lands within noise of the
label-shuffle floor (Table~\ref{tab:label-shuffle}) on both datasets,
while $H_0$ alone is statistically indistinguishable from the fused
$H_0{+}H_1$ default -- essentially all the recoverable signal is in $H_0$
(connected-component birth/death, i.e.\ local density) and $H_1$ (loop
structure) contributes nothing measurable beyond it on either design
distribution. The Betti curve does not replicate this collapse: BC's
$H_1$-only arm ($0.260$ T / $0.510$ NT) is worse than its $H_0$-only arm
($0.237$ T / $0.294$ NT) but nowhere near the $\approx\!0.9$--$1.0$
label-shuffle floor that PI's $H_1$-only arm collapses to, on either
dataset. A feature-count curve therefore retains some usable signal in the
loop structure that a persistence-image raster, at this calibration, does
not -- though BC's $H_0$-only number is itself well above PI's, so this is
not a claim that BC's $H_1$ channel is more informative in absolute terms.

\subsubsection{CoordConv}
\label{ssec:branch-coordconv}

We also augment each persistence image with two coordinate channels, on the
reasoning that calibration is what makes a pixel position mean something and the
network should be handed that rather than have to recover it. The ablation does
not support the reasoning, so the headline model does not use them.

During the processing step each $(D, 64, 64)$ persistence-image
stack is augmented with two coordinate channels -- fixed $x$- and
$y$-coordinate grids linearly spaced over $[-1, 1]$ -- following the CoordConv
construction REF [Liu et al., CoordConv 2018], so that the convolutional stack has access
to absolute pixel position rather than having to infer it. This ensures that the work
done during the calibration step is not lost: each pixel has a fixed,
shared meaning across the dataset and CoordConv lets the network exploit this
directly instead of relearning it.

\begin{table}[tbp]
\centering
\small
\begin{tabular}{lcc}
\toprule
Arm & Thomas $\mathcal{L}$ & Nested-Thomas $\mathcal{L}$ \\
\midrule
Full model (CoordConv on) & $0.124 \pm 0.007$ & $0.196 \pm 0.006$ \\
CoordConv removed & $0.126 \pm 0.005$ & $0.194 \pm 0.002$ \\
\bottomrule
\end{tabular}
\caption{CoordConv ablation: the shared-weight encoder of
Section~\ref{ssec:encoder-fusion} with and without the two coordinate
channels, at the default $q=0.95$ calibration, $n=5$ seeds per arm on both
design distributions.}
\label{tab:coordconv-ablation}
\end{table}

CoordConv removal is within seed noise of the full model on both datasets --
if anything marginally better on Thomas and on nested-Thomas. This does not
support the claim that CoordConv is necessary to exploit the calibrated pixel
grid; at the current calibration ($q=0.95$, Section~\ref{ssec:calibration})
the convolutional stack appears to recover comparable information without
explicit coordinate channels.

\subsubsection{Encoder and fusion mode}
\label{ssec:exp-ablations}

We also cross encoder mode against fusion mode. Flat concatenation under a
shared encoder is both the simplest arm and the only one that trains reliably,
which is what Section~\ref{ssec:encoder-fusion} adopts -- though with a single
scale the encoder-sharing axis of this cross does not arise at all.

The architecture sweep was run on the persistence-image arm at the default $q=0.95$
calibration; Section~\ref{ssec:exp-summaries} confirms PI is the arm worth
ablating.

The sweep crosses \emph{encoder mode} (shared-weight vs.\
independent-per-$k$) against \emph{fusion mode} (concatenation, average-pool,
or a small convolutional-flatten fusion head), five seeds per combination, on
Thomas.

\begin{table}[tbp]
\centering
\small
\begin{tabular}{llcc}
\toprule
Encoder mode & Fusion & $\mathcal{L}$ (mean $\pm$ sd) & Failed seeds \\
\midrule
shared      & concat        & $0.122 \pm 0.007$ & $0/5$ \\
shared      & conv-avg      & $0.431 \pm 0.425^\dagger$ & $2/5$ \\
shared      & conv-flatten  & $0.589 \pm 0.425^\dagger$ & $3/5$ \\
independent & concat        & $0.129 \pm 0.006$ & $0/5$ \\
independent & conv-avg      & $0.133 \pm 0.015$ & $0/5$ \\
independent & conv-flatten  & $0.141 \pm 0.016$ & $0/5$ \\
\bottomrule
\end{tabular}
\caption{Full encoder-mode $\times$ fusion-mode cross, $n=5$ seeds per
combination, on Thomas at the default PI calibration. ``Failed'' seeds are
those that collapse to the $\mathcal{L}\approx0.9$ design-distribution-mean
floor of Table~\ref{tab:label-shuffle} rather than converging.
$^\dagger$ mean is dominated by the failed seeds and is not a stable estimate
of the arm's typical performance; see per-seed values in the text.}
\label{tab:encoder-fusion-sweep}
\end{table}

\fig{Per-seed loss scatter across the six configurations -- the collapse
pattern below is the interesting part and this table's means alone
understate it.}

Restricted to concatenation, shared and independent encoders match to within
seed noise ($0.122\pm0.007$ vs.\ $0.129\pm0.006$); since sharing weights across
scales reduces the parameter count (roughly $K$-fold) and removes a design
axis, we adopt the shared encoder going forward. 

Both convolutional-fusion variants are unstable, but \emph{only when paired with the shared encoder}:
shared/conv-avg and shared/conv-flatten each collapse to the
$\mathcal{L}\approx0.9$ floor on 2/5 and 3/5 seeds respectively (individual
per-seed losses: $\{0.128, 0.120, 0.902, 0.115, 0.892\}$ and $\{0.903, 0.118,
0.903, 0.130, 0.892\}$), while the identical fusion heads under an
independent-per-$k$ encoder converge cleanly on all five seeds every time
($0.133\pm0.015$, $0.141\pm0.016$) without matching plain concatenation.

Flat concatenation is therefore preferred not only for its simplicity but
because -- together with the shared-encoder counterpart -- it is the only
combination that trains reliably and reaches the best loss. The same two concatenation arms (shared and independent encoder) were
re-run on nested-Thomas to check the choice transfers across design
distributions:

\begin{table}[tbp]
\centering
\small
\begin{tabular}{lcc}
\toprule
Encoder mode & Thomas $\mathcal{L}$ & Nested-Thomas $\mathcal{L}$ \\
\midrule
shared, concat      & $0.124 \pm 0.007$ & $0.196 \pm 0.006$ \\
independent, concat & $0.133 \pm 0.005$ & $0.228 \pm 0.029$ \\
\bottomrule
\end{tabular}
\caption{Cross-dataset confirmation of the shared-vs-independent encoder
choice, $n=5$ seeds per arm.}
\label{tab:encoder-confirm}
\end{table}

On nested-Thomas the shared encoder's advantage over the independent one is
larger than on Thomas and the independent arm is noisier
($\pm0.029$ vs.\ $\pm0.005$), so the shared-encoder choice seems to matter even
more on the more difficult five-parameter design distribution, reinforcing 
the choice made above.

\subsubsection{Alternative vectorizations}
\label{ssec:exp-summaries}

We also implemented three alternative vectorizations of the same diagrams --
landscapes, silhouettes, and Betti/Euler curves. None of them improves on the
persistence image, on either process or either encoder path, so the persistence
image is the summary carried through the paper.
Figure~\ref{fig:vectorization_schematic} gives a visual overview of the
vectorizations, computed from the same persistence diagram; we describe each in
turn and report their respective performance in Table~\ref{tab:main-comparison}.

\begin{figure}[tbp]
    \centering
    \includegraphics[width=\textwidth]{vectorization_schematic.pdf}
    \caption{One point cloud $\rightarrow$ its $H_0/H_1$ persistence diagrams
    $\rightarrow$ the vectorizations compared in
    Section~\ref{ssec:exp-summaries}: persistence-image raster, landscape
    stack, silhouette curve, Betti curve, and persistence-statistics summary.
    Top row: the tight-cluster design ($\kappa=50$, $\mu=4$, $\sigma=0.01$);
    bottom row: the near-Poisson design ($\kappa=4$, $\mu=5$, $\sigma=0.5$),
    the same two extremes of Table~\ref{tab:sim-validation}. Statistics-panel
    bars are min-max normalized per statistic, with raw values printed above
    each bar, since count, total persistence, and entropy live on
    incomparable scales.}
    \label{fig:vectorization_schematic}
\end{figure}

\paragraph{Landscape.} The $K$ layers $\lambda_1(t) \geq \lambda_2(t) \geq
\dots \geq \lambda_K(t)$ are the order statistics of the tent family
$\{\Lambda_i(t)\}_i$ (Definition~\ref{def:tent}) at each grid point $t$:
$\lambda_j(t)$ is the $j$-th largest tent value. $K$ is chosen per diagram
sample by a coverage rule -- the smallest $K$ for which layer $K+1$ is
negligible for 99\% of training diagrams, capped at 16 -- giving a $(K,G)$
raster per homology dimension.

\paragraph{Silhouette.} Rather than keep every order statistic, the
silhouette collapses the same tents into one persistence-weighted average,
    \[
    \phi^{(p)}(t) = \frac{\sum_i |d_i-b_i|^p\, \Lambda_i(t)}
    {\sum_i |d_i-b_i|^p},
    \]
a single $(1,G)$ curve, with the amplitude exponent $p$ fixed at $1$ here.

\paragraph{Betti/Euler curve.} The Betti curve $\beta_k(r) = \#\{i : b_i
\leq r < d_i\}$ counts $k$-features alive at scale $r$ -- a feature count
rather than a tent reduction, but built from the same persistence pairs.
The Euler curve $\chi(r) = \beta_0(r) - \beta_1(r)$ is a signed count nested
inside the two Betti channels, not additional information.

For a discussion of each vectorization's calibration (grid ranges, $K$-selection, shared vs.\ per-dimension
axes) refer to Appendix~\ref{app:calibration-sweeps}; for per-stage tensor
shapes refer to Appendix~\ref{app:tensor-shapes}.

\begin{table}[tbp]
\centering
\footnotesize
\begin{tabular}{lrrcccc}
\toprule
 & & & \multicolumn{2}{c}{$\mathcal{L}$, native encoder} & \multicolumn{2}{c}{$\mathcal{L}$, MLP encoder} \\
\cmidrule(lr){4-5}\cmidrule(lr){6-7}
Vec. & dim/channel & channels & T & NT & T & NT \\
\midrule
PI  & $64\times64=4096$ & 6  & $0.126 \pm 0.008$ (10) & $0.196 \pm 0.006$ (5) & $0.154 \pm 0.012$ (5) & $0.261 \pm 0.011$ (5) \\
LS  & $16\times128=2048$ & 6  & $0.132 \pm 0.008$ (3)  & $0.250 \pm 0.005$ (3) & $0.167 \pm 0.019$ (5) & $0.284 \pm 0.009$ (5) \\
SIL & $128$              & 12 & $0.165 \pm 0.013$ (3)  & $0.271 \pm 0.012$ (3) & $0.177 \pm 0.013$ (5) & $0.297 \pm 0.004$ (5) \\
BC  & $128$              & 6  & $0.196 \pm 0.014$ (3)  & $0.290 \pm 0.005$ (3) & $0.176 \pm 0.009$ (5) & $0.266 \pm 0.003$ (5) \\
\bottomrule
\end{tabular}
\caption{Vectorization methods comparison. ``Native'' denotes the methods'
purpose-built encoder (CNN for PI/LS, 1-D CoordConv for SIL/BC); ``MLP'' is
the architecture-agnostic \texttt{FlattenMLPEncoder} applied to the same
tensor. All native-path arms use the $q=0.95$ calibration matched to the
persistence image rather than each
vectorization's own tuned default from Appendix~\ref{app:alt-calibration},
so the comparison is calibration-fair;
BC's own appendix-preferred $q=0.99$ scores slightly better
($0.179\pm0.012$ T / $0.284\pm0.011$ NT) but is not used here for
consistency. SIL is fixed at $p=1$ throughout (the SIL-1 arm of
Appendix~\ref{app:calib-silhouette}) for both the native and MLP path. T =
Thomas, NT = nested-Thomas; seed count $n$ in parentheses.}
\label{tab:main-comparison}
\end{table}

The persistence image is the strongest arm on both design distributions and on
both encoder paths, and the margin survives the substitution of the
architecture-agnostic MLP encoder for each vectorization's own, so the ranking
is not an artifact of the encoder each summary happens to be paired with.


% #############################################################################
% 8. DISCUSSION AND LIMITATIONS
% #############################################################################
% NOTE -- PURPOSE: pre-empt the obvious objections. Short.
% #############################################################################
\section{Discussion and Limitations}
\label{sec:discussion}

\emph{Notes, not prose. Each bullet is one paragraph to write.}
\begin{itemize}
    \item \textbf{Negative results.} Three alternative vectorizations were
    implemented and none improved on persistence images
    (Section~\ref{ssec:exp-summaries}). Say plainly what this does and does not
    license -- it is evidence about these summaries under this architecture and
    this design distribution, not a general claim about vectorizations. The
    same qualification covers the other branches of
    Section~\ref{ssec:branches}: multi-scale fusion, $H_1$, and coordinate
    channels each fail to earn their place \emph{here}.
    \item \textbf{Where the simplification costs something.} Fused scales beat
    $DTM_5$ on nested Thomas ($0.196$ vs.\ $0.205$,
    Section~\ref{ssec:branch-multiscale}), so the single-scale model is a trade
    on the two-level process and not a free saving; say so rather than let the
    Thomas result stand for both.
    \item \textbf{Calibration inherited across summaries.} The per-diagram
    coverage rule was inherited unchanged by the 1-D summaries for
    comparability, which may not be optimal for each individually.
    \item \textbf{Compute cost.} \tab{Feature-extraction time per cloud by
    vectorization, training time, and amortized inference time per cloud
    against minimum contrast per cloud. The last comparison strongly favours
    the method and is currently unmentioned.}
    \item \textbf{The design distribution is a prior.} The design distribution
    acts as a prior, so the trained model approximates a posterior summary
    under it, and its ranges must therefore cover the intended application
    domain. This is the single most important limitation of any amortized
    estimator and deserves a full paragraph, not a trailing sentence.
    \item \textbf{Computational cost.} Persistence computation scales badly in
    $n$; this is why the point budget of Section~\ref{ssec:design} is capped at
    800. State the cost explicitly and compare against the cost of minimum
    contrast.
    \item \textbf{Single realization.} The estimator sees one realization;
    classical methods often assume the same, but say so explicitly.
    \item \textbf{Homogeneity and stationarity.} Everything here assumes
    homogeneous, stationary, isotropic processes.
\end{itemize}

% #############################################################################
% 9. FUTURE WORK
% #############################################################################
% NOTE -- PURPOSE: REQUIRED for the progression report; for the paper this can
% be folded into Section 8. Keep it concrete and sequenced -- this is the
% section the panel reads to judge whether there is a viable thesis here.
% #############################################################################
\section{Future Work}
\label{sec:futurework}

\emph{Notes, not prose. Sequenced as written.}
\begin{itemize}
    \item \textbf{Immediate.} Complete Section~\ref{ssec:exp-performance}: the
    headline confirmation run, the error floor, the grid surface, the
    extrapolation set, and the Mat\'ern cluster runs. Give this a timeline --
    it is the nearest-term deliverable.
    \item \textbf{Near term.} The encoder question: attention over scales, and
    set-based encoders operating on diagrams directly rather than on rasters.
    \item \textbf{Medium term.} Inhomogeneous / non-stationary processes, where
    classical closed forms are least available and the topological approach has
    the most to gain -- this is arguably where the thesis-scale contribution
    lies.
    \item \textbf{Medium term.} Real data. Galaxy catalogues if the running
    example of Section~\ref{sec:intro} is cosmological; name a specific
    dataset.
    \item \textbf{Longer term.} Uncertainty quantification -- the model
    currently emits a point estimate with no interval, which limits its use as
    a drop-in replacement for likelihood-based methods.
\end{itemize}

% #############################################################################
% APPENDIX
% #############################################################################
% NOTE -- PURPOSE:
%   A. every calibration sweep in the paper, one subsection per vectorization,
%      in a common format -- the persistence-image sweep included.
%   B. per-stage tensor shapes for all vectorizations. The material exists in
%      the implementation notes and needs pasting in.
% #############################################################################
\appendix

\section{Calibration Sweeps}
\label{app:calibration-sweeps}

\label{app:alt-calibration}% retained: referenced from Section~\ref{ssec:exp-summaries}

Section~\ref{ssec:calibration} discusses the calibration for the persistence image's birth/persistence
axes by per-diagram coverage quantiles fit on the training split alone. The same
$(q, \alpha)$ scheme -- coverage $q$ and padding $\alpha$ over per-diagram
$b^{\min}, b^{\max}, p^{\max}$ -- carries over unchanged to the 1-D grid
$[t_{\min}, T]$ underlying landscapes, silhouettes, and Betti/Euler curves, since
all three are built from the same tent family $\{\Lambda_i(t)\}_i$
(Definition~\ref{def:tent}) rather than the 2-D birth-persistence surface. This
appendix reports, per vectorization -- the persistence image included -- the
calibration sweeps run at the multi-scale architecture of
Section~\ref{ssec:branch-multiscale}: a CoordConv-augmented
CNN encoder with shared weights across DTM scales $k \in \{5,10,15\}$ followed
by concatenation fusion, and a
two-layer MLP head. Every sweep below therefore predates the simplifications of
Section~\ref{sec:method} and was run under the branch configuration, not under
the headline model.

\subsection{Persistence Images}
\label{app:calib-pi}

We quantify the effect
of the calibration parameters introduced in Section~\ref{ssec:calibration}:
coverage $q$, padding $\alpha$, kernel width $\sigma_{\text{pixels}}$, and
raster resolution $G$. Table~\ref{tab:calibration-sweep} reports the
sweep, each arm varying a single parameter around the default
calibration while holding architecture and the remaining three parameters
fixed.

\fig{$\mathcal{L}$ against $q$, both design distributions, with $q=1.00$ marked
as the naive endpoint. One-parameter family, so a line plot reads better than a
table and makes the dataset disagreement below visible at a glance.}

\begin{table}[tbp]
\centering
\small
\begin{tabular}{llcc}
\toprule
Swept parameter & Value & Thomas $\mathcal{L}$ & Nested-Thomas $\mathcal{L}$ \\
\midrule
Coverage $q$
  & 0.90            & $0.125 \pm 0.008$ & $0.195 \pm 0.004$ \\
  & 0.95 (default)  & $0.126 \pm 0.006$ & $0.441 \pm 0.341^\dagger$ \\
  & 0.99            & $0.122 \pm 0.006$ & $0.203 \pm 0.009$ \\
  & 1.00 (naive)    & $0.120 \pm 0.005$ & $0.210 \pm 0.004$ \\
\addlinespace
Padding $\alpha$
  & 1.00            & $0.123 \pm 0.007$ & $0.201 \pm 0.008$ \\
  & 1.05 (default)  & $0.126 \pm 0.007$ & $0.199 \pm 0.007$ \\
  & 1.10            & $0.125 \pm 0.006$ & $0.202 \pm 0.009$ \\
\addlinespace
Kernel width $\sigma_{\text{pixels}}$
  & 0.5             & $0.124 \pm 0.007$ & $0.192 \pm 0.008$ \\
  & 1.0             & $0.123 \pm 0.006$ & $0.202 \pm 0.013$ \\
  & 2.0             & $0.124 \pm 0.009$ & $0.457 \pm 0.329^\dagger$ \\
\addlinespace
Resolution $G$
  & 32              & $0.122 \pm 0.004$ & $0.209 \pm 0.011$ \\
  & 64 (default)    & $0.125 \pm 0.009$ & $0.199 \pm 0.007$ \\
  & 128             & $0.635 \pm 0.352^\dagger$ & $0.684 \pm 0.339^\dagger$ \\
\bottomrule
\end{tabular}
% Channel bookkeeping (not restated in the caption -- see the (B,n_k,D,G,G)
% shape in Section~\ref{ssec:branch-multiscale}): n_k=3 DTM scales x D=2 homology
% dimensions (H0+H1) = 6 PI channels total, each G x G.
\caption{Calibration sweep for the persistence-image (PI) vectorization,
one-at-a-time around the default calibration ($q=0.95$, $\alpha=1.05$,
$\sigma_{\text{pixels}}=0.75$, $G=64$; Section~\ref{ssec:calibration}), for
both the Thomas and nested-Thomas design distributions
(Section~\ref{sec:simulator}). Diagrams use the default $DTM_{\{5,10,15\}}$,
$H_0{+}H_1$ pipeline and shared-weight architecture of
Section~\ref{ssec:encoder-fusion}. Cells report $\mathcal{L}$
(Section~\ref{ssec:eval-protocol}, mean $\pm$ sd over 3 seeds).
$^\dagger$ marks configurations that diverge on at least one of the
three seeds (per-seed $\mathcal{L} \approx 0.9$--$1.0$, i.e.\ no better than
predicting the design-distribution mean); their mean is dominated by the
diverged seed(s) and should not be read as a stable estimate.}
\label{tab:calibration-sweep}
\end{table}

Excluding the $^\dagger$ rows, resolution $G=128$ is the only arm that is
unambiguously bad -- $G\in\{32,64\}$, $\alpha$, and $\sigma_{\text{pixels}}$
all vary within one seed-standard-deviation of the default across both
datasets. Coverage $q$ is the one parameter with a genuine process-dependent
disagreement: the Thomas optimum moves toward the naive endpoint $q=1.00$
($0.120$ vs.\ $0.126$ at the default), while nested-Thomas prefers tighter
coverage at $q=0.90$ ($0.195$) and gets markedly worse as $q\to1$. We choose
to keep the shared default $q=0.95$.

The $^\dagger$ rows are the one behaviour the calibration does not explain.
$G=128$ fails on at least one of three seeds on \emph{both} design
distributions, which looks like an optimization failure at the increased raster
size rather than a calibration effect, and the same reading applies to the
single diverged $\sigma_{\text{pixels}}=2.0$ seed on nested Thomas; neither is a
statement about the calibrated range.

\subsection{Persistence Landscapes}
\label{app:calib-landscape}

For homology dimension $h$ and DTM scale $k$, the landscape layers
$\lambda_1(t) \geq \lambda_2(t) \geq \dots \geq \lambda_K(t)$ are the order
statistics of $\{\Lambda_i(t)\}_i$ at each grid point $t$ (Section~\ref{ssec:exp-summaries}),
giving a $(K, G)$ raster per $(h, k)$ that is fed to the encoder as a
non-square image (height $K$, width $G$) -- the only structural difference from
the $G \times G$ persistence image. $K$ is not hand-picked: a coverage rule
selects, per $(h,k)$, the smallest $K$ such that
    \[
    \|\lambda_{K+1}\|_\infty \;\leq\; \tau \cdot \|\lambda_1\|_\infty
    \]
for at least a fraction $c$ of training diagrams, capped at $K_{\max}$; the
per-dimension answers are then reconciled to their shared maximum so that the
$H_0$ and $H_1$ channels can be stacked into one raster. Three calibration axes
are therefore specific to this vectorization and are not present in the
persistence-image sweep: the grid resolution $G$, the rule's own relative
threshold $\tau$ (fixed at $0.01$ throughout every other experiment in this
paper), and $K$ itself, which we also sweep as a fixed override against letting
the rule choose it. Table~\ref{tab:calibration-sweep-landscape} reports all
three, one-at-a-time, holding $c = 0.99$, $K_{\max} = 16$, and the diagram-grid
coverage $q = 0.95$ fixed except where swept.

\begin{table}[tbp]
\centering
\small
\begin{tabular}{llcc}
\toprule
Swept parameter & Value & Thomas $\mathcal{L}$ & Nested-Thomas $\mathcal{L}$ \\
\midrule
Grid resolution $G$
  & 64             & $0.135 \pm 0.015$ & $0.249 \pm 0.003$ \\
  & 128 (default)  & $0.137 \pm 0.008$ & $0.240 \pm 0.003$ \\
  & 256            & $0.141 \pm 0.014$ & $0.246 \pm 0.002$ \\
\addlinespace
$K$-rule threshold $\tau$
  & 0.5\%            & $0.134 \pm 0.008$ & $0.245 \pm 0.007$ \\
  & 1\% (default)    & $0.135 \pm 0.010$ & $0.248 \pm 0.006$ \\
  & 2\%              & $0.146 \pm 0.026$ & $0.243 \pm 0.006$ \\
\addlinespace
Fixed $K$ (vs.\ data-derived)
  & 4                                  & $0.174 \pm 0.013$ & $0.372 \pm 0.023$ \\
  & 8 (budget-matched, $G=512$)$^\ddagger$ & $0.154 \pm 0.011$ & $0.312 \pm 0.007$ \\
  & 16                                 & $0.132 \pm 0.012$ & $0.244 \pm 0.006$ \\
  & data-derived (default)             & $0.133 \pm 0.007$ & $0.245 \pm 0.006$ \\
\bottomrule
\end{tabular}
\caption{Calibration sweep for the persistence-landscape (LS) vectorization,
one-at-a-time around the native-budget default ($G=128$, $\tau=0.01$, $K$
data-derived), on the \emph{params} task for both the Thomas and nested-Thomas
design distributions (Section~\ref{sec:simulator}). Diagrams use the default
$DTM_{\{5,10,15\}}$, $H_0{+}H_1$ pipeline and shared-weight architecture of
Section~\ref{ssec:encoder-fusion} -- the same one used for the
persistence-image sweep of Table~\ref{tab:calibration-sweep} -- with two
landscape channels per scale (six total). The diagram-grid coverage $q =
0.95$ (Section~\ref{ssec:calibration}) is fixed across every arm;
$K$-selection uses coverage $c=0.99$ and cap $K_{\max}=16$, both fixed except
in the third group. Cells report $\mathcal{L}$
(Section~\ref{ssec:eval-protocol}, mean $\pm$ sd over 3 seeds). $^\ddagger$
the $G=512$, $K=8$ arm matches the $4096$-d persistence-image budget of
Table~\ref{tab:calibration-sweep} exactly; the other three rows of that group
hold $G=128$ fixed and vary only $K$, so this row is not raster-size-matched
to them.}
\label{tab:calibration-sweep-landscape}
\end{table}

\textbf{The coverage rule resolved to $K=16$ -- the cap -- on every arm and
seed underlying the first two groups}, so the $G$ and $\tau$ sweeps report
robustness at a $K$ already pinned by $K_{\max}$, not by $\tau$ itself; their
near-flat $\mathcal{L}$ (differences mostly within one seed-sd) is consistent
with sweeping a quantity ($G$, $\tau$) that is not presently the binding
constraint on $K$. The fixed-$K$ group is the informative one for this rule:
the ``data-derived'' row lands within seed noise of the explicit $K=16$ row
on both datasets, exactly as expected since both resolve to the identical
($G=128$, $K=16$) configuration, while $K=4$ and $K=8$ are both markedly and
monotonically worse. Within the range tested, more landscape layers help and
$K_{\max}$ -- not $\tau$ -- is the binding constraint; distinguishing the
rule's own threshold from its cap would require re-running with a larger
$K_{\max}$.
\subsection{Silhouettes}
\label{app:calib-silhouette}

The silhouette collapses the same tent family $\{\Lambda_i(t)\}_i$ into the
single persistence-weighted curve $\phi^{(p)}(t)$ of
Section~\ref{ssec:exp-summaries}, sampled on the same coverage-quantile
grid $[t_{\min}, T]$ (Section~\ref{ssec:calibration}) at resolution $G$; both
the normalized curve and its unnormalized numerator are kept as separate
channels (the denominator discards the feature count, which the numerator
preserves), giving $2 \times 2 \times 3 = 12$ 1-D channels ($H_0{+}H_1$,
DTM $k \in \{5,10,15\}$) fed to a per-scale 1-D CoordConv-style encoder under
the same shared-weight, concat-fusion architecture as
Table~\ref{tab:calibration-sweep}. Unlike $K$ for landscapes, the amplitude
exponent $p$ has no data-derived rule to calibrate: Section~\ref{ssec:exp-summaries}
treats it as an experimental arm compared against $\{0,2,3\}$ elsewhere, not a
calibration choice, so this appendix fixes $p=1$ (the SIL-1 arm) and reports
only the grid-resolution sweep, holding $q=0.95$ and $\alpha=1.05$
(Section~\ref{ssec:calibration}) fixed at the values selected for the
persistence image.

\begin{table}[tbp]
\centering
\small
\begin{tabular}{llcc}
\toprule
Swept parameter & Value & Thomas $\mathcal{L}$ & Nested-Thomas $\mathcal{L}$ \\
\midrule
Grid resolution $G$
  & 64             & $0.169 \pm 0.008$ & $0.268 \pm 0.006$ \\
  & 128 (default)  & $0.165 \pm 0.011$ & $0.271 \pm 0.009$ \\
  & 256            & $0.170 \pm 0.013$ & $0.273 \pm 0.012$ \\
\bottomrule
\end{tabular}
\caption{Grid-resolution calibration sweep for the silhouette (SIL-1, $p=1$)
vectorization, on the \emph{params} task for both the Thomas and
nested-Thomas design distributions (Section~\ref{sec:simulator}). Diagrams
use the default $DTM_{\{5,10,15\}}$, $H_0{+}H_1$ pipeline and shared-weight
architecture of Section~\ref{ssec:encoder-fusion} -- the same one used for
the persistence-image sweep of Table~\ref{tab:calibration-sweep}. The
diagram-grid coverage $q=0.95$ and padding $\alpha=1.05$
(Section~\ref{ssec:calibration}) are fixed across every arm at the values
selected for the persistence image. Cells report $\mathcal{L}$
(Section~\ref{ssec:eval-protocol}, mean $\pm$ sd over 3 seeds).}
\label{tab:calibration-sweep-silhouette}
\end{table}

Resolution has little effect in either direction: every arm is within
roughly one seed-sd of every other on both design distributions, and no
configuration shows the divergence behavior marked $^\dagger$ in
Table~\ref{tab:calibration-sweep}. $G=64$ is a reasonable default going
forward -- it costs a quarter of $G=256$'s raster and is not measurably worse
on either dataset.
\subsection{Betti and Euler Curves}
\label{app:calib-betti}

The Betti curve $\beta_k(r)$ and Euler curve $\chi(r) = \beta_0(r) -
\beta_1(r)$ of Section~\ref{ssec:exp-summaries} are sampled on the same
coverage-quantile grid as the other 1-D vectorizations, giving a
$(\text{grid\_size},)$ curve per homology dimension and DTM scale. The runs
below use $\chi$ disabled (two Betti channels stacked, $H_0{+}H_1$, no
derived Euler channel), so the channel count matches SIL-1 and the landscape
sweep -- 2 dims $\times$ 3 scales = 6 channels -- rather than the
$H_0{+}H_1{+}\chi$ default used elsewhere in the paper. A Betti curve has no
layer-count axis, so there is no $K$-analogue to calibrate here, and its
values are already bounded (small alive-feature counts, not raw pixel
magnitudes), so grid resolution $G$ and calibration coverage $q$ are the two
axes swept, one at a time, at the shared multi-scale architecture of
Table~\ref{tab:calibration-sweep}.

\begin{table}[tbp]
\centering
\small
\begin{tabular}{llcc}
\toprule
Swept parameter & Value & Thomas $\mathcal{L}$ & Nested-Thomas $\mathcal{L}$ \\
\midrule
Grid resolution $G$
  & 64             & $0.183 \pm 0.009$ & $0.284 \pm 0.006$ \\
  & 128 (default)  & $0.200 \pm 0.004$ & $0.289 \pm 0.002$ \\
  & 256            & $0.196 \pm 0.012$ & $0.293 \pm 0.008$ \\
\addlinespace
Coverage $q$
  & 0.95 (default) & $0.190 \pm 0.008$ & $0.291 \pm 0.003$ \\
  & 0.99           & $0.178 \pm 0.009$ & $0.285 \pm 0.007$ \\
  & 1.00 (naive)   & $0.187 \pm 0.009$ & $0.344 \pm 0.002$ \\
\bottomrule
\end{tabular}
\caption{Calibration sweep for the Betti-curve (BC) vectorization, one-at-a-time
around the native-budget default (grid\_size $=128$, $q=0.95$), on the
\emph{params} task for both the Thomas and nested-Thomas design distributions
(Section~\ref{sec:simulator}). Diagrams use the default $DTM_{\{5,10,15\}}$,
$H_0{+}H_1$ pipeline (Euler channel disabled, see text) and shared-weight
architecture of Section~\ref{ssec:encoder-fusion}. Padding is fixed at
$\alpha=1.05$ throughout -- the value selected for the persistence image --
rather than this vectorization's own code default of $\alpha=1.1$. Cells
report $\mathcal{L}$ (Section~\ref{ssec:eval-protocol}, mean $\pm$ sd over
3 seeds).}
\label{tab:calibration-sweep-betti}
\end{table}

$G=64$ is again at least as good as $G=128$ on both design distributions --
the same pattern as SIL-1 above, consistent with a curve-shaped vectorization
having little to gain from a finer grid once the underlying feature count is
small. Coverage behaves differently than it does for the persistence image:
$q=0.99$ is the best arm on both datasets, and $q=1.00$ (the naive,
uncalibrated range) is clearly worse on nested-Thomas -- a $\sim$21\% increase
in $\mathcal{L}$ over $q=0.99$, with a seed-sd tight enough ($\pm0.002$) that
this is a real effect and not a diverged seed. Unlike the $^\dagger$ rows of
Table~\ref{tab:calibration-sweep}, no arm here shows per-seed divergence to
the design-distribution mean; $q=1.00$ is simply a worse calibrated range, not
a failure mode.


\section{Tensor Shapes}
\label{app:tensor-shapes}

The persistence-image path of Section~\ref{ssec:encoder-fusion} runs
$(M_i, 2)$ birth--death pairs $\rightarrow$ one $64 \times 64$ raster per
diagram $\rightarrow (B, 1, 64, 64) \rightarrow (B, E)$ out of the encoder
$\rightarrow (B, E+1)$ at the head input; the multi-scale variant of
Section~\ref{ssec:branch-multiscale} carries the extra $n_k$ axis folded into
the batch dimension, giving $(B\!\cdot\!n_k, D{+}2, 64, 64) \rightarrow
(B, n_k, E) \rightarrow (B, n_k\!\cdot\!E + 1)$ with the two coordinate
channels of Section~\ref{ssec:branch-coordconv} included.

\tab{One stage-by-stage shape table per vectorization, from $(M_i, 2)$
birth--death pairs through to the head input, in the notation of
Section~\ref{ssec:encoder-fusion} ($B$, $n_k$, $D$, $G$, $E$). The landscape
$(K,G)$ raster, silhouette two-channel curve, Betti curve and statistics vector
paths need the same treatment. Flag in the landscape row that the raster is
non-square, which is why the coordinate-channel builder needs independent height
and width ramps.}


\end{document}